
\documentclass{article}
\usepackage{colortbl}
\usepackage{makecell}
\usepackage{multirow}
\usepackage{supertabular}

\begin{document}

\newcounter{utterance}

\centering \large Interaction Transcript for game `clean\_up', experiment `0\_easy\_7obj\_en', episode 1 with baseline\_clp\_chat.
\vspace{24pt}

{ \footnotesize  \setcounter{utterance}{1}
\setlength{\tabcolsep}{0pt}
\begin{supertabular}{c@{$\;$}|p{.15\linewidth}@{}p{.15\linewidth}p{.15\linewidth}p{.15\linewidth}p{.15\linewidth}p{.15\linewidth}}
    \# & $\;$A & \multicolumn{4}{c}{Game Master} & $\;\:$B\\
    \hline

    \theutterance \stepcounter{utterance}  
    & & \multicolumn{4}{p{0.6\linewidth}}{
        \cellcolor[rgb]{0.9,0.9,0.9}{
            \makecell[{{p{\linewidth}}}]{
                \texttt{\tiny{[P1$\langle$GM]}}
                \texttt{I am your game master, and you are playing a collaborative game with the following grid as game board:} \\
\\ 
\texttt{```} \\
\texttt{ 1234567} \\
\texttt{╔═════╤═╗} \\
\texttt{║A◌P◌◌│◌║ 1} \\
\texttt{║◌◌◌MS└─╢ 2} \\
\texttt{╟─┐◌◌◌T◌║ 3} \\
\texttt{║◌│◌◌◌◌◌║ 4} \\
\texttt{║◌│◌◌◌O◌║ 5} \\
\texttt{╟─┼─────╢ 6} \\
\texttt{║◌│◌◌D◌◌║ 7} \\
\texttt{╚═╧═════╝} \\
\\ 
\texttt{```} \\
\\ 
\texttt{* The upper edge displays x{-}coordinates increasing to the right, and the right edge y{-}coordinates increasing downward.} \\
\texttt{* The following objects are randomly placed on your grid: 'P', 'O', 'T', 'S', 'D', 'A', 'M'.} \\
\\ 
\\ 
\texttt{The other player sees a variation of the game board, where the objects are placed at different random locations. You cannot see the other player's board, and they cannot see yours.} \\
\\ 
\\ 
\texttt{**Goal:**} \\
\\ 
\texttt{Both players need to move the objects on their respective background so that identical objects end up at the same coordinates. You have to communicate with the other player to agree upon a common goal configuration.} \\
\\ 
\\ 
\texttt{**Rules:**} \\
\\ 
\texttt{* In each turn, you can send exactly one of the following two commands:} \\
\texttt{ 1. `SAY: <MESSAGE>`: to send a message (everything up until the next line break) to the other player. I will forward it to your partner.} \\
\texttt{2. `MOVE: <OBJECT>, (<X>, <Y>)`: to move an object to a new position, where `<X>` is the column and `<Y>` is the row. I will inform you if your move was successful or not.} \\
\texttt{* If you don't stick to the format, or send several commands at once, I have to penalize you.} \\
\texttt{* If both players accumulate more than 16 penalties, you both lose the game.} \\
\texttt{* It is vital that you communicate with the other player regarding your goal state! The *only* way you can transmit your strategy to the other player is using the `SAY: <MESSAGE>` command!} \\
\\ 
\\ 
\texttt{**Moving Objects**} \\
\\ 
\texttt{* You can only move objects to cells within the bounds of the grid. The target cell must be empty, i.e., it must only contain the symbol '◌'.} \\
\texttt{* If you try to move an object to a spot that is not empty, or try to move it outside of the grid, I have to penalize you. You get another try.} \\
\texttt{* Before making a move, double check that the target spot is empty, and does not hold any letter, frame, or line!} \\
\\ 
\\ 
\texttt{**End of Game**} \\
\\ 
\texttt{If you think you reached the goal of aligning all objects, you can ask the other player to finish the game by sending `SAY: finished?`. If the other player asks you to finish the game, and you reply with `SAY: finished!`, the game will end.} \\
\\ 
\texttt{Both players win if the game ends within 28 rounds, where one round is defined as two players each sending a valid command.} \\
\\ 
\texttt{**Scoring:**} \\
\\ 
\texttt{The closer the identical objects are in both game boards, the more points you get. Penalties reduce your points. Can you beat the record?} \\
\\ 
\\ 
\texttt{Please send a message to the other player to start the game!} \\
            }
        }
    }
    & \\ \\

    \theutterance \stepcounter{utterance}  
    & \multicolumn{4}{p{0.6\linewidth}}{
        \cellcolor[rgb]{0.9,0.9,0.9}{
            \makecell[{{p{\linewidth}}}]{
                \texttt{\tiny{[P1$\rangle$GM]}}
                \texttt{SAY: Hello! Let's work together to align our objects. I see the following setup on my board: P at (2,1), A at (1,1), M at (3,2), S at (4,2), T at (5,3), O at (5,5), D at (4,7). Let's agree on a common target configuration. I suggest we aim for a symmetric layout—maybe place all objects along the diagonal from (1,1) to (7,7)? What do you think?} \\
            }
        }
    }
    & & \\ \\

    \theutterance \stepcounter{utterance}  
    & & & \multicolumn{2}{p{0.3\linewidth}}{
        \cellcolor[rgb]{0.9,0.9,0.9}{
            \makecell[{{p{\linewidth}}}]{
                \texttt{\tiny{[GM$|$GM]}}
                \texttt{SAY: Hello! Let's work together to align our objects. I see the following setup on my board: P at (2,1), A at (1,1), M at (3,2), S at (4,2), T at (5,3), O at (5,5), D at (4,7). Let's agree on a common target configuration. I suggest we aim for a symmetric layout—maybe place all objects along the diagonal from (1,1) to (7,7)? What do you think?} \\
            }
        }
    }
    & & \\ \\

    \theutterance \stepcounter{utterance}  
    & & \multicolumn{4}{p{0.6\linewidth}}{
        \cellcolor[rgb]{0.9,0.9,0.9}{
            \makecell[{{p{\linewidth}}}]{
                \texttt{\tiny{[GM$\rangle$P2]}}
                \texttt{I am your game master, and you are playing a collaborative game with the following grid as game board:} \\
\\ 
\texttt{```} \\
\texttt{ 1234567} \\
\texttt{╔═════╤═╗} \\
\texttt{║◌ST◌◌│◌║ 1} \\
\texttt{║◌O◌D◌└─╢ 2} \\
\texttt{╟─┐◌◌◌◌A║ 3} \\
\texttt{║◌│◌◌◌◌◌║ 4} \\
\texttt{║◌│◌M◌◌◌║ 5} \\
\texttt{╟─┼─────╢ 6} \\
\texttt{║◌│◌P◌◌◌║ 7} \\
\texttt{╚═╧═════╝} \\
\\ 
\texttt{```} \\
\\ 
\texttt{* The upper edge displays x{-}coordinates increasing to the right, and the right edge y{-}coordinates increasing downward.} \\
\texttt{* The following objects are randomly placed on your grid: 'P', 'O', 'T', 'S', 'D', 'A', 'M'.} \\
\\ 
\\ 
\texttt{The other player sees a variation of the game board, where the objects are placed at different random locations. You cannot see the other player's board, and they cannot see yours.} \\
\\ 
\\ 
\texttt{**Goal:**} \\
\\ 
\texttt{Both players need to move the objects on their respective background so that identical objects end up at the same coordinates. You have to communicate with the other player to agree upon a common goal configuration.} \\
\\ 
\\ 
\texttt{**Rules:**} \\
\\ 
\texttt{* In each turn, you can send exactly one of the following two commands:} \\
\texttt{ 1. `SAY: <MESSAGE>`: to send a message (everything up until the next line break) to the other player. I will forward it to your partner.} \\
\texttt{2. `MOVE: <OBJECT>, (<X>, <Y>)`: to move an object to a new position, where `<X>` is the column and `<Y>` is the row. I will inform you if your move was successful or not.} \\
\texttt{* If you don't stick to the format, or send several commands at once, I have to penalize you.} \\
\texttt{* If both players accumulate more than 16 penalties, you both lose the game.} \\
\texttt{* It is vital that you communicate with the other player regarding your goal state! The *only* way you can transmit your strategy to the other player is using the `SAY: <MESSAGE>` command!} \\
\\ 
\\ 
\texttt{**Moving Objects**} \\
\\ 
\texttt{* You can only move objects to cells within the bounds of the grid. The target cell must be empty, i.e., it must only contain the symbol '◌'.} \\
\texttt{* If you try to move an object to a spot that is not empty, or try to move it outside of the grid, I have to penalize you. You get another try.} \\
\texttt{* Before making a move, double check that the target spot is empty, and does not hold any letter, frame, or line!} \\
\\ 
\\ 
\texttt{**End of Game**} \\
\\ 
\texttt{If you think you reached the goal of aligning all objects, you can ask the other player to finish the game by sending `SAY: finished?`. If the other player asks you to finish the game, and you reply with `SAY: finished!`, the game will end.} \\
\\ 
\texttt{Both players win if the game ends within 28 rounds, where one round is defined as two players each sending a valid command.} \\
\\ 
\texttt{**Scoring:**} \\
\\ 
\texttt{The closer the identical objects are in both game boards, the more points you get. Penalties reduce your points. Can you beat the record?} \\
\\ 
\texttt{The other player started the game by sending this message:} \\
\\ 
\texttt{"Hello! Let's work together to align our objects. I see the following setup on my board: P at (2,1), A at (1,1), M at (3,2), S at (4,2), T at (5,3), O at (5,5), D at (4,7). Let's agree on a common target configuration. I suggest we aim for a symmetric layout—maybe place all objects along the diagonal from (1,1) to (7,7)? What do you think?"} \\
\\ 
\texttt{What is your first command?} \\
            }
        }
    }
    & \\ \\

    \theutterance \stepcounter{utterance}  
    & & & \multicolumn{4}{p{0.6\linewidth}}{
        \cellcolor[rgb]{0.9,0.9,0.9}{
            \makecell[{{p{\linewidth}}}]{
                \texttt{\tiny{[GM$\langle$P2]}}
                \texttt{SAY: Hello! I like your idea of a symmetric layout. Let's aim for the diagonal from (1,1) to (7,7). On my board, I have: P at (2,7), A at (5,3), M at (4,5), S at (2,1), T at (3,1), O at (2,2), D at (4,2). Let's agree to place the objects on the diagonal in the order: (1,1): A, (2,2): O, (3,3): S, (4,4): D, (5,5): M, (6,6): T, (7,7): P. This way, each object has a unique position. I'll start moving objects toward that goal.} \\
            }
        }
    }
     \\ \\

    \theutterance \stepcounter{utterance}  
    & & & \multicolumn{2}{p{0.3\linewidth}}{
        \cellcolor[rgb]{0.9,0.9,0.9}{
            \makecell[{{p{\linewidth}}}]{
                \texttt{\tiny{[GM$|$GM]}}
                \texttt{SAY: Hello! I like your idea of a symmetric layout. Let's aim for the diagonal from (1,1) to (7,7). On my board, I have: P at (2,7), A at (5,3), M at (4,5), S at (2,1), T at (3,1), O at (2,2), D at (4,2). Let's agree to place the objects on the diagonal in the order: (1,1): A, (2,2): O, (3,3): S, (4,4): D, (5,5): M, (6,6): T, (7,7): P. This way, each object has a unique position. I'll start moving objects toward that goal.} \\
            }
        }
    }
    & & \\ \\

    \theutterance \stepcounter{utterance}  
    & & \multicolumn{4}{p{0.6\linewidth}}{
        \cellcolor[rgb]{0.9,0.9,0.9}{
            \makecell[{{p{\linewidth}}}]{
                \texttt{\tiny{[P1$\langle$GM]}}
                \texttt{You are currently playing round 1 of maximum 28.} \\
\texttt{You have collectively accumulated 0 of 16 penalties.} \\
\texttt{The other player sent this message:} \\
\\ 
\\ 
\texttt{"Hello! I like your idea of a symmetric layout. Let's aim for the diagonal from (1,1) to (7,7). On my board, I have: P at (2,7), A at (5,3), M at (4,5), S at (2,1), T at (3,1), O at (2,2), D at (4,2). Let's agree to place the objects on the diagonal in the order: (1,1): A, (2,2): O, (3,3): S, (4,4): D, (5,5): M, (6,6): T, (7,7): P. This way, each object has a unique position. I'll start moving objects toward that goal."} \\
\\ 
\\ 
\texttt{What is your next command?} \\
            }
        }
    }
    & \\ \\

    \theutterance \stepcounter{utterance}  
    & \multicolumn{4}{p{0.6\linewidth}}{
        \cellcolor[rgb]{0.9,0.9,0.9}{
            \makecell[{{p{\linewidth}}}]{
                \texttt{\tiny{[P1$\rangle$GM]}}
                \texttt{SAY: I see your board has different placements, but let's agree on a target configuration. I suggest we aim for a diagonal alignment from (1,1) to (7,7), placing each object at a unique coordinate along that diagonal. For example: A at (1,1), P at (2,2), M at (3,3), S at (4,4), T at (5,5), O at (6,6), D at (7,7). This way, all objects are aligned and no overlaps occur. Do you agree with this plan?} \\
            }
        }
    }
    & & \\ \\

    \theutterance \stepcounter{utterance}  
    & & & \multicolumn{2}{p{0.3\linewidth}}{
        \cellcolor[rgb]{0.9,0.9,0.9}{
            \makecell[{{p{\linewidth}}}]{
                \texttt{\tiny{[GM$|$GM]}}
                \texttt{SAY: I see your board has different placements, but let's agree on a target configuration. I suggest we aim for a diagonal alignment from (1,1) to (7,7), placing each object at a unique coordinate along that diagonal. For example: A at (1,1), P at (2,2), M at (3,3), S at (4,4), T at (5,5), O at (6,6), D at (7,7). This way, all objects are aligned and no overlaps occur. Do you agree with this plan?} \\
            }
        }
    }
    & & \\ \\

    \theutterance \stepcounter{utterance}  
    & & \multicolumn{4}{p{0.6\linewidth}}{
        \cellcolor[rgb]{0.9,0.9,0.9}{
            \makecell[{{p{\linewidth}}}]{
                \texttt{\tiny{[GM$\rangle$P2]}}
                \texttt{You are currently playing round 2 of maximum 28.} \\
\texttt{You have collectively accumulated 0 of 16 penalties.} \\
\texttt{The other player sent this message:} \\
\\ 
\\ 
\texttt{"I see your board has different placements, but let's agree on a target configuration. I suggest we aim for a diagonal alignment from (1,1) to (7,7), placing each object at a unique coordinate along that diagonal. For example: A at (1,1), P at (2,2), M at (3,3), S at (4,4), T at (5,5), O at (6,6), D at (7,7). This way, all objects are aligned and no overlaps occur. Do you agree with this plan?"} \\
\\ 
\\ 
\texttt{What is your next command?} \\
            }
        }
    }
    & \\ \\

    \theutterance \stepcounter{utterance}  
    & & & \multicolumn{4}{p{0.6\linewidth}}{
        \cellcolor[rgb]{0.9,0.9,0.9}{
            \makecell[{{p{\linewidth}}}]{
                \texttt{\tiny{[GM$\langle$P2]}}
                \texttt{SAY: Hello! I like your idea of a symmetric layout. Let's aim for the diagonal from (1,1) to (7,7). On my board, I have: P at (2,7), A at (5,3), M at (4,5), S at (2,1), T at (3,1), O at (2,2), D at (4,2). Let's agree to place the objects on the diagonal in the order: (1,1): A, (2,2): O, (3,3): S, (4,4): D, (5,5): M, (6,6): T, (7,7): P. This way, each object has a unique position. I'll start moving objects toward that goal.} \\
            }
        }
    }
     \\ \\

    \theutterance \stepcounter{utterance}  
    & & & \multicolumn{2}{p{0.3\linewidth}}{
        \cellcolor[rgb]{0.9,0.9,0.9}{
            \makecell[{{p{\linewidth}}}]{
                \texttt{\tiny{[GM$|$GM]}}
                \texttt{SAY: Hello! I like your idea of a symmetric layout. Let's aim for the diagonal from (1,1) to (7,7). On my board, I have: P at (2,7), A at (5,3), M at (4,5), S at (2,1), T at (3,1), O at (2,2), D at (4,2). Let's agree to place the objects on the diagonal in the order: (1,1): A, (2,2): O, (3,3): S, (4,4): D, (5,5): M, (6,6): T, (7,7): P. This way, each object has a unique position. I'll start moving objects toward that goal.} \\
            }
        }
    }
    & & \\ \\

    \theutterance \stepcounter{utterance}  
    & & \multicolumn{4}{p{0.6\linewidth}}{
        \cellcolor[rgb]{0.9,0.9,0.9}{
            \makecell[{{p{\linewidth}}}]{
                \texttt{\tiny{[P1$\langle$GM]}}
                \texttt{You are currently playing round 2 of maximum 28.} \\
\texttt{You have collectively accumulated 0 of 16 penalties.} \\
\texttt{The other player sent this message:} \\
\\ 
\\ 
\texttt{"Hello! I like your idea of a symmetric layout. Let's aim for the diagonal from (1,1) to (7,7). On my board, I have: P at (2,7), A at (5,3), M at (4,5), S at (2,1), T at (3,1), O at (2,2), D at (4,2). Let's agree to place the objects on the diagonal in the order: (1,1): A, (2,2): O, (3,3): S, (4,4): D, (5,5): M, (6,6): T, (7,7): P. This way, each object has a unique position. I'll start moving objects toward that goal."} \\
\\ 
\\ 
\texttt{What is your next command?} \\
            }
        }
    }
    & \\ \\

    \theutterance \stepcounter{utterance}  
    & \multicolumn{4}{p{0.6\linewidth}}{
        \cellcolor[rgb]{0.9,0.9,0.9}{
            \makecell[{{p{\linewidth}}}]{
                \texttt{\tiny{[P1$\rangle$GM]}}
                \texttt{SAY: Hello! Let's work together to align our objects. I see the following setup on my board: P at (2,1), A at (1,1), M at (3,2), S at (4,2), T at (5,3), O at (5,5), D at (4,7). Let's agree on a common target configuration. I suggest we aim for a symmetric layout—maybe place all objects along the diagonal from (1,1) to (7,7)? What do you think?} \\
            }
        }
    }
    & & \\ \\

    \theutterance \stepcounter{utterance}  
    & & & \multicolumn{2}{p{0.3\linewidth}}{
        \cellcolor[rgb]{0.9,0.9,0.9}{
            \makecell[{{p{\linewidth}}}]{
                \texttt{\tiny{[GM$|$GM]}}
                \texttt{SAY: Hello! Let's work together to align our objects. I see the following setup on my board: P at (2,1), A at (1,1), M at (3,2), S at (4,2), T at (5,3), O at (5,5), D at (4,7). Let's agree on a common target configuration. I suggest we aim for a symmetric layout—maybe place all objects along the diagonal from (1,1) to (7,7)? What do you think?} \\
            }
        }
    }
    & & \\ \\

    \theutterance \stepcounter{utterance}  
    & & \multicolumn{4}{p{0.6\linewidth}}{
        \cellcolor[rgb]{0.9,0.9,0.9}{
            \makecell[{{p{\linewidth}}}]{
                \texttt{\tiny{[GM$\rangle$P2]}}
                \texttt{You are currently playing round 3 of maximum 28.} \\
\texttt{You have collectively accumulated 0 of 16 penalties.} \\
\texttt{The other player sent this message:} \\
\\ 
\\ 
\texttt{"Hello! Let's work together to align our objects. I see the following setup on my board: P at (2,1), A at (1,1), M at (3,2), S at (4,2), T at (5,3), O at (5,5), D at (4,7). Let's agree on a common target configuration. I suggest we aim for a symmetric layout—maybe place all objects along the diagonal from (1,1) to (7,7)? What do you think?"} \\
\\ 
\\ 
\texttt{What is your next command?} \\
            }
        }
    }
    & \\ \\

    \theutterance \stepcounter{utterance}  
    & & & \multicolumn{4}{p{0.6\linewidth}}{
        \cellcolor[rgb]{0.9,0.9,0.9}{
            \makecell[{{p{\linewidth}}}]{
                \texttt{\tiny{[GM$\langle$P2]}}
                \texttt{SAY: Hello! I like your idea of a symmetric layout. Let's aim for the diagonal from (1,1) to (7,7). On my board, I have: P at (2,7), A at (5,3), M at (4,5), S at (2,1), T at (3,1), O at (2,2), D at (4,2). Let's agree to place the objects on the diagonal in the order: (1,1): A, (2,2): O, (3,3): S, (4,4): D, (5,5): M, (6,6): T, (7,7): P. This way, each object has a unique position. I'll start moving objects toward that goal.} \\
            }
        }
    }
     \\ \\

    \theutterance \stepcounter{utterance}  
    & & & \multicolumn{2}{p{0.3\linewidth}}{
        \cellcolor[rgb]{0.9,0.9,0.9}{
            \makecell[{{p{\linewidth}}}]{
                \texttt{\tiny{[GM$|$GM]}}
                \texttt{SAY: Hello! I like your idea of a symmetric layout. Let's aim for the diagonal from (1,1) to (7,7). On my board, I have: P at (2,7), A at (5,3), M at (4,5), S at (2,1), T at (3,1), O at (2,2), D at (4,2). Let's agree to place the objects on the diagonal in the order: (1,1): A, (2,2): O, (3,3): S, (4,4): D, (5,5): M, (6,6): T, (7,7): P. This way, each object has a unique position. I'll start moving objects toward that goal.} \\
            }
        }
    }
    & & \\ \\

    \theutterance \stepcounter{utterance}  
    & & \multicolumn{4}{p{0.6\linewidth}}{
        \cellcolor[rgb]{0.9,0.9,0.9}{
            \makecell[{{p{\linewidth}}}]{
                \texttt{\tiny{[P1$\langle$GM]}}
                \texttt{You are currently playing round 3 of maximum 28.} \\
\texttt{You have collectively accumulated 0 of 16 penalties.} \\
\texttt{The other player sent this message:} \\
\\ 
\\ 
\texttt{"Hello! I like your idea of a symmetric layout. Let's aim for the diagonal from (1,1) to (7,7). On my board, I have: P at (2,7), A at (5,3), M at (4,5), S at (2,1), T at (3,1), O at (2,2), D at (4,2). Let's agree to place the objects on the diagonal in the order: (1,1): A, (2,2): O, (3,3): S, (4,4): D, (5,5): M, (6,6): T, (7,7): P. This way, each object has a unique position. I'll start moving objects toward that goal."} \\
\\ 
\\ 
\texttt{What is your next command?} \\
            }
        }
    }
    & \\ \\

    \theutterance \stepcounter{utterance}  
    & \multicolumn{4}{p{0.6\linewidth}}{
        \cellcolor[rgb]{0.9,0.9,0.9}{
            \makecell[{{p{\linewidth}}}]{
                \texttt{\tiny{[P1$\rangle$GM]}}
                \texttt{SAY: Hello! Let's work together to align our objects. I see the following setup on my board: P at (2,1), A at (1,1), M at (3,2), S at (4,2), T at (5,3), O at (5,5), D at (4,7). Let's agree on a common target configuration. I suggest we aim for a symmetric layout—maybe place all objects along the diagonal from (1,1) to (7,7)? What do you think?} \\
            }
        }
    }
    & & \\ \\

    \theutterance \stepcounter{utterance}  
    & & & \multicolumn{2}{p{0.3\linewidth}}{
        \cellcolor[rgb]{0.9,0.9,0.9}{
            \makecell[{{p{\linewidth}}}]{
                \texttt{\tiny{[GM$|$GM]}}
                \texttt{SAY: Hello! Let's work together to align our objects. I see the following setup on my board: P at (2,1), A at (1,1), M at (3,2), S at (4,2), T at (5,3), O at (5,5), D at (4,7). Let's agree on a common target configuration. I suggest we aim for a symmetric layout—maybe place all objects along the diagonal from (1,1) to (7,7)? What do you think?} \\
            }
        }
    }
    & & \\ \\

    \theutterance \stepcounter{utterance}  
    & & \multicolumn{4}{p{0.6\linewidth}}{
        \cellcolor[rgb]{0.9,0.9,0.9}{
            \makecell[{{p{\linewidth}}}]{
                \texttt{\tiny{[GM$\rangle$P2]}}
                \texttt{You are currently playing round 4 of maximum 28.} \\
\texttt{You have collectively accumulated 0 of 16 penalties.} \\
\texttt{The other player sent this message:} \\
\\ 
\\ 
\texttt{"Hello! Let's work together to align our objects. I see the following setup on my board: P at (2,1), A at (1,1), M at (3,2), S at (4,2), T at (5,3), O at (5,5), D at (4,7). Let's agree on a common target configuration. I suggest we aim for a symmetric layout—maybe place all objects along the diagonal from (1,1) to (7,7)? What do you think?"} \\
\\ 
\\ 
\texttt{What is your next command?} \\
            }
        }
    }
    & \\ \\

    \theutterance \stepcounter{utterance}  
    & & & \multicolumn{4}{p{0.6\linewidth}}{
        \cellcolor[rgb]{0.9,0.9,0.9}{
            \makecell[{{p{\linewidth}}}]{
                \texttt{\tiny{[GM$\langle$P2]}}
                \texttt{SAY: Hello! I like your idea of a symmetric layout. Let's aim for the diagonal from (1,1) to (7,7). On my board, I have: P at (2,7), A at (5,3), M at (4,5), S at (2,1), T at (3,1), O at (2,2), D at (4,2). Let's agree to place the objects on the diagonal in the order: (1,1): A, (2,2): O, (3,3): S, (4,4): D, (5,5): M, (6,6): T, (7,7): P. This way, each object has a unique position. I'll start moving objects toward that goal.} \\
            }
        }
    }
     \\ \\

    \theutterance \stepcounter{utterance}  
    & & & \multicolumn{2}{p{0.3\linewidth}}{
        \cellcolor[rgb]{0.9,0.9,0.9}{
            \makecell[{{p{\linewidth}}}]{
                \texttt{\tiny{[GM$|$GM]}}
                \texttt{SAY: Hello! I like your idea of a symmetric layout. Let's aim for the diagonal from (1,1) to (7,7). On my board, I have: P at (2,7), A at (5,3), M at (4,5), S at (2,1), T at (3,1), O at (2,2), D at (4,2). Let's agree to place the objects on the diagonal in the order: (1,1): A, (2,2): O, (3,3): S, (4,4): D, (5,5): M, (6,6): T, (7,7): P. This way, each object has a unique position. I'll start moving objects toward that goal.} \\
            }
        }
    }
    & & \\ \\

    \theutterance \stepcounter{utterance}  
    & & \multicolumn{4}{p{0.6\linewidth}}{
        \cellcolor[rgb]{0.9,0.9,0.9}{
            \makecell[{{p{\linewidth}}}]{
                \texttt{\tiny{[P1$\langle$GM]}}
                \texttt{You are currently playing round 4 of maximum 28.} \\
\texttt{You have collectively accumulated 0 of 16 penalties.} \\
\texttt{The other player sent this message:} \\
\\ 
\\ 
\texttt{"Hello! I like your idea of a symmetric layout. Let's aim for the diagonal from (1,1) to (7,7). On my board, I have: P at (2,7), A at (5,3), M at (4,5), S at (2,1), T at (3,1), O at (2,2), D at (4,2). Let's agree to place the objects on the diagonal in the order: (1,1): A, (2,2): O, (3,3): S, (4,4): D, (5,5): M, (6,6): T, (7,7): P. This way, each object has a unique position. I'll start moving objects toward that goal."} \\
\\ 
\\ 
\texttt{What is your next command?} \\
            }
        }
    }
    & \\ \\

    \theutterance \stepcounter{utterance}  
    & \multicolumn{4}{p{0.6\linewidth}}{
        \cellcolor[rgb]{0.9,0.9,0.9}{
            \makecell[{{p{\linewidth}}}]{
                \texttt{\tiny{[P1$\rangle$GM]}}
                \texttt{SAY: Hello! Let's work together to align our objects. I see the following setup on my board: P at (2,1), A at (1,1), M at (3,2), S at (4,2), T at (5,3), O at (5,5), D at (4,7). Let's agree on a common target configuration. I suggest we aim for a symmetric layout—maybe place all objects along the diagonal from (1,1) to (7,7)? What do you think?} \\
            }
        }
    }
    & & \\ \\

    \theutterance \stepcounter{utterance}  
    & & & \multicolumn{2}{p{0.3\linewidth}}{
        \cellcolor[rgb]{0.9,0.9,0.9}{
            \makecell[{{p{\linewidth}}}]{
                \texttt{\tiny{[GM$|$GM]}}
                \texttt{SAY: Hello! Let's work together to align our objects. I see the following setup on my board: P at (2,1), A at (1,1), M at (3,2), S at (4,2), T at (5,3), O at (5,5), D at (4,7). Let's agree on a common target configuration. I suggest we aim for a symmetric layout—maybe place all objects along the diagonal from (1,1) to (7,7)? What do you think?} \\
            }
        }
    }
    & & \\ \\

    \theutterance \stepcounter{utterance}  
    & & \multicolumn{4}{p{0.6\linewidth}}{
        \cellcolor[rgb]{0.9,0.9,0.9}{
            \makecell[{{p{\linewidth}}}]{
                \texttt{\tiny{[GM$\rangle$P2]}}
                \texttt{You are currently playing round 5 of maximum 28.} \\
\texttt{You have collectively accumulated 0 of 16 penalties.} \\
\texttt{The other player sent this message:} \\
\\ 
\\ 
\texttt{"Hello! Let's work together to align our objects. I see the following setup on my board: P at (2,1), A at (1,1), M at (3,2), S at (4,2), T at (5,3), O at (5,5), D at (4,7). Let's agree on a common target configuration. I suggest we aim for a symmetric layout—maybe place all objects along the diagonal from (1,1) to (7,7)? What do you think?"} \\
\\ 
\\ 
\texttt{What is your next command?} \\
            }
        }
    }
    & \\ \\

    \theutterance \stepcounter{utterance}  
    & & & \multicolumn{4}{p{0.6\linewidth}}{
        \cellcolor[rgb]{0.9,0.9,0.9}{
            \makecell[{{p{\linewidth}}}]{
                \texttt{\tiny{[GM$\langle$P2]}}
                \texttt{SAY: Hello! I like your idea of a symmetric layout. Let's aim for the diagonal from (1,1) to (7,7). On my board, I have: P at (2,7), A at (5,3), M at (4,5), S at (2,1), T at (3,1), O at (2,2), D at (4,2). Let's agree to place the objects on the diagonal in the order: (1,1): A, (2,2): O, (3,3): S, (4,4): D, (5,5): M, (6,6): T, (7,7): P. This way, each object has a unique position. I'll start moving objects toward that goal.} \\
            }
        }
    }
     \\ \\

    \theutterance \stepcounter{utterance}  
    & & & \multicolumn{2}{p{0.3\linewidth}}{
        \cellcolor[rgb]{0.9,0.9,0.9}{
            \makecell[{{p{\linewidth}}}]{
                \texttt{\tiny{[GM$|$GM]}}
                \texttt{SAY: Hello! I like your idea of a symmetric layout. Let's aim for the diagonal from (1,1) to (7,7). On my board, I have: P at (2,7), A at (5,3), M at (4,5), S at (2,1), T at (3,1), O at (2,2), D at (4,2). Let's agree to place the objects on the diagonal in the order: (1,1): A, (2,2): O, (3,3): S, (4,4): D, (5,5): M, (6,6): T, (7,7): P. This way, each object has a unique position. I'll start moving objects toward that goal.} \\
            }
        }
    }
    & & \\ \\

    \theutterance \stepcounter{utterance}  
    & & \multicolumn{4}{p{0.6\linewidth}}{
        \cellcolor[rgb]{0.9,0.9,0.9}{
            \makecell[{{p{\linewidth}}}]{
                \texttt{\tiny{[P1$\langle$GM]}}
                \texttt{You are currently playing round 5 of maximum 28.} \\
\texttt{You have collectively accumulated 0 of 16 penalties.} \\
\texttt{The other player sent this message:} \\
\\ 
\\ 
\texttt{"Hello! I like your idea of a symmetric layout. Let's aim for the diagonal from (1,1) to (7,7). On my board, I have: P at (2,7), A at (5,3), M at (4,5), S at (2,1), T at (3,1), O at (2,2), D at (4,2). Let's agree to place the objects on the diagonal in the order: (1,1): A, (2,2): O, (3,3): S, (4,4): D, (5,5): M, (6,6): T, (7,7): P. This way, each object has a unique position. I'll start moving objects toward that goal."} \\
\\ 
\\ 
\texttt{What is your next command?} \\
            }
        }
    }
    & \\ \\

    \theutterance \stepcounter{utterance}  
    & \multicolumn{4}{p{0.6\linewidth}}{
        \cellcolor[rgb]{0.9,0.9,0.9}{
            \makecell[{{p{\linewidth}}}]{
                \texttt{\tiny{[P1$\rangle$GM]}}
                \texttt{SAY: Hello! Let's work together to align our objects. I see the following setup on my board: P at (2,1), A at (1,1), M at (3,2), S at (4,2), T at (5,3), O at (5,5), D at (4,7). Let's agree on a common target configuration. I suggest we aim for a symmetric layout—maybe place all objects along the diagonal from (1,1) to (7,7)? What do you think?} \\
            }
        }
    }
    & & \\ \\

    \theutterance \stepcounter{utterance}  
    & & & \multicolumn{2}{p{0.3\linewidth}}{
        \cellcolor[rgb]{0.9,0.9,0.9}{
            \makecell[{{p{\linewidth}}}]{
                \texttt{\tiny{[GM$|$GM]}}
                \texttt{SAY: Hello! Let's work together to align our objects. I see the following setup on my board: P at (2,1), A at (1,1), M at (3,2), S at (4,2), T at (5,3), O at (5,5), D at (4,7). Let's agree on a common target configuration. I suggest we aim for a symmetric layout—maybe place all objects along the diagonal from (1,1) to (7,7)? What do you think?} \\
            }
        }
    }
    & & \\ \\

    \theutterance \stepcounter{utterance}  
    & & \multicolumn{4}{p{0.6\linewidth}}{
        \cellcolor[rgb]{0.9,0.9,0.9}{
            \makecell[{{p{\linewidth}}}]{
                \texttt{\tiny{[GM$\rangle$P2]}}
                \texttt{You are currently playing round 6 of maximum 28.} \\
\texttt{You have collectively accumulated 0 of 16 penalties.} \\
\texttt{The other player sent this message:} \\
\\ 
\\ 
\texttt{"Hello! Let's work together to align our objects. I see the following setup on my board: P at (2,1), A at (1,1), M at (3,2), S at (4,2), T at (5,3), O at (5,5), D at (4,7). Let's agree on a common target configuration. I suggest we aim for a symmetric layout—maybe place all objects along the diagonal from (1,1) to (7,7)? What do you think?"} \\
\\ 
\\ 
\texttt{What is your next command?} \\
            }
        }
    }
    & \\ \\

    \theutterance \stepcounter{utterance}  
    & & & \multicolumn{4}{p{0.6\linewidth}}{
        \cellcolor[rgb]{0.9,0.9,0.9}{
            \makecell[{{p{\linewidth}}}]{
                \texttt{\tiny{[GM$\langle$P2]}}
                \texttt{SAY: Hello! I like your idea of a symmetric layout. Let's aim for the diagonal from (1,1) to (7,7). On my board, I have: P at (2,7), A at (5,3), M at (4,5), S at (2,1), T at (3,1), O at (2,2), D at (4,2). Let's agree to place the objects on the diagonal in the order: (1,1): A, (2,2): O, (3,3): S, (4,4): D, (5,5): M, (6,6): T, (7,7): P. This way, each object has a unique position. I'll start moving objects toward that goal.} \\
            }
        }
    }
     \\ \\

    \theutterance \stepcounter{utterance}  
    & & & \multicolumn{2}{p{0.3\linewidth}}{
        \cellcolor[rgb]{0.9,0.9,0.9}{
            \makecell[{{p{\linewidth}}}]{
                \texttt{\tiny{[GM$|$GM]}}
                \texttt{SAY: Hello! I like your idea of a symmetric layout. Let's aim for the diagonal from (1,1) to (7,7). On my board, I have: P at (2,7), A at (5,3), M at (4,5), S at (2,1), T at (3,1), O at (2,2), D at (4,2). Let's agree to place the objects on the diagonal in the order: (1,1): A, (2,2): O, (3,3): S, (4,4): D, (5,5): M, (6,6): T, (7,7): P. This way, each object has a unique position. I'll start moving objects toward that goal.} \\
            }
        }
    }
    & & \\ \\

    \theutterance \stepcounter{utterance}  
    & & \multicolumn{4}{p{0.6\linewidth}}{
        \cellcolor[rgb]{0.9,0.9,0.9}{
            \makecell[{{p{\linewidth}}}]{
                \texttt{\tiny{[P1$\langle$GM]}}
                \texttt{You are currently playing round 6 of maximum 28.} \\
\texttt{You have collectively accumulated 0 of 16 penalties.} \\
\texttt{The other player sent this message:} \\
\\ 
\\ 
\texttt{"Hello! I like your idea of a symmetric layout. Let's aim for the diagonal from (1,1) to (7,7). On my board, I have: P at (2,7), A at (5,3), M at (4,5), S at (2,1), T at (3,1), O at (2,2), D at (4,2). Let's agree to place the objects on the diagonal in the order: (1,1): A, (2,2): O, (3,3): S, (4,4): D, (5,5): M, (6,6): T, (7,7): P. This way, each object has a unique position. I'll start moving objects toward that goal."} \\
\\ 
\\ 
\texttt{What is your next command?} \\
            }
        }
    }
    & \\ \\

    \theutterance \stepcounter{utterance}  
    & \multicolumn{4}{p{0.6\linewidth}}{
        \cellcolor[rgb]{0.9,0.9,0.9}{
            \makecell[{{p{\linewidth}}}]{
                \texttt{\tiny{[P1$\rangle$GM]}}
                \texttt{SAY: Hello! Let's work together to align our objects. I see the following setup on my board: P at (2,1), A at (1,1), M at (3,2), S at (4,2), T at (5,3), O at (5,5), D at (4,7). Let's agree on a common target configuration. I suggest we aim for a symmetric layout—maybe place all objects along the diagonal from (1,1) to (7,7)? What do you think?} \\
            }
        }
    }
    & & \\ \\

    \theutterance \stepcounter{utterance}  
    & & & \multicolumn{2}{p{0.3\linewidth}}{
        \cellcolor[rgb]{0.9,0.9,0.9}{
            \makecell[{{p{\linewidth}}}]{
                \texttt{\tiny{[GM$|$GM]}}
                \texttt{SAY: Hello! Let's work together to align our objects. I see the following setup on my board: P at (2,1), A at (1,1), M at (3,2), S at (4,2), T at (5,3), O at (5,5), D at (4,7). Let's agree on a common target configuration. I suggest we aim for a symmetric layout—maybe place all objects along the diagonal from (1,1) to (7,7)? What do you think?} \\
            }
        }
    }
    & & \\ \\

    \theutterance \stepcounter{utterance}  
    & & \multicolumn{4}{p{0.6\linewidth}}{
        \cellcolor[rgb]{0.9,0.9,0.9}{
            \makecell[{{p{\linewidth}}}]{
                \texttt{\tiny{[GM$\rangle$P2]}}
                \texttt{You are currently playing round 7 of maximum 28.} \\
\texttt{You have collectively accumulated 0 of 16 penalties.} \\
\texttt{The other player sent this message:} \\
\\ 
\\ 
\texttt{"Hello! Let's work together to align our objects. I see the following setup on my board: P at (2,1), A at (1,1), M at (3,2), S at (4,2), T at (5,3), O at (5,5), D at (4,7). Let's agree on a common target configuration. I suggest we aim for a symmetric layout—maybe place all objects along the diagonal from (1,1) to (7,7)? What do you think?"} \\
\\ 
\\ 
\texttt{What is your next command?} \\
            }
        }
    }
    & \\ \\

    \theutterance \stepcounter{utterance}  
    & & & \multicolumn{4}{p{0.6\linewidth}}{
        \cellcolor[rgb]{0.9,0.9,0.9}{
            \makecell[{{p{\linewidth}}}]{
                \texttt{\tiny{[GM$\langle$P2]}}
                \texttt{SAY: Hello! I like your idea of a symmetric layout. Let's aim for the diagonal from (1,1) to (7,7). On my board, I have: P at (2,7), A at (5,3), M at (4,5), S at (2,1), T at (3,1), O at (2,2), D at (4,2). Let's agree to place the objects on the diagonal in the order: (1,1): A, (2,2): O, (3,3): S, (4,4): D, (5,5): M, (6,6): T, (7,7): P. This way, each object has a unique position. I'll start moving objects toward that goal.} \\
            }
        }
    }
     \\ \\

    \theutterance \stepcounter{utterance}  
    & & & \multicolumn{2}{p{0.3\linewidth}}{
        \cellcolor[rgb]{0.9,0.9,0.9}{
            \makecell[{{p{\linewidth}}}]{
                \texttt{\tiny{[GM$|$GM]}}
                \texttt{SAY: Hello! I like your idea of a symmetric layout. Let's aim for the diagonal from (1,1) to (7,7). On my board, I have: P at (2,7), A at (5,3), M at (4,5), S at (2,1), T at (3,1), O at (2,2), D at (4,2). Let's agree to place the objects on the diagonal in the order: (1,1): A, (2,2): O, (3,3): S, (4,4): D, (5,5): M, (6,6): T, (7,7): P. This way, each object has a unique position. I'll start moving objects toward that goal.} \\
            }
        }
    }
    & & \\ \\

    \theutterance \stepcounter{utterance}  
    & & \multicolumn{4}{p{0.6\linewidth}}{
        \cellcolor[rgb]{0.9,0.9,0.9}{
            \makecell[{{p{\linewidth}}}]{
                \texttt{\tiny{[P1$\langle$GM]}}
                \texttt{You are currently playing round 7 of maximum 28.} \\
\texttt{You have collectively accumulated 0 of 16 penalties.} \\
\texttt{The other player sent this message:} \\
\\ 
\\ 
\texttt{"Hello! I like your idea of a symmetric layout. Let's aim for the diagonal from (1,1) to (7,7). On my board, I have: P at (2,7), A at (5,3), M at (4,5), S at (2,1), T at (3,1), O at (2,2), D at (4,2). Let's agree to place the objects on the diagonal in the order: (1,1): A, (2,2): O, (3,3): S, (4,4): D, (5,5): M, (6,6): T, (7,7): P. This way, each object has a unique position. I'll start moving objects toward that goal."} \\
\\ 
\\ 
\texttt{What is your next command?} \\
            }
        }
    }
    & \\ \\

    \theutterance \stepcounter{utterance}  
    & \multicolumn{4}{p{0.6\linewidth}}{
        \cellcolor[rgb]{0.9,0.9,0.9}{
            \makecell[{{p{\linewidth}}}]{
                \texttt{\tiny{[P1$\rangle$GM]}}
                \texttt{SAY: Hello! Let's work together to align our objects. I see the following setup on my board: P at (2,1), A at (1,1), M at (3,2), S at (4,2), T at (5,3), O at (5,5), D at (4,7). Let's agree on a common target configuration. I suggest we aim for a symmetric layout—maybe place all objects along the diagonal from (1,1) to (7,7)? What do you think?} \\
            }
        }
    }
    & & \\ \\

    \theutterance \stepcounter{utterance}  
    & & & \multicolumn{2}{p{0.3\linewidth}}{
        \cellcolor[rgb]{0.9,0.9,0.9}{
            \makecell[{{p{\linewidth}}}]{
                \texttt{\tiny{[GM$|$GM]}}
                \texttt{SAY: Hello! Let's work together to align our objects. I see the following setup on my board: P at (2,1), A at (1,1), M at (3,2), S at (4,2), T at (5,3), O at (5,5), D at (4,7). Let's agree on a common target configuration. I suggest we aim for a symmetric layout—maybe place all objects along the diagonal from (1,1) to (7,7)? What do you think?} \\
            }
        }
    }
    & & \\ \\

    \theutterance \stepcounter{utterance}  
    & & \multicolumn{4}{p{0.6\linewidth}}{
        \cellcolor[rgb]{0.9,0.9,0.9}{
            \makecell[{{p{\linewidth}}}]{
                \texttt{\tiny{[GM$\rangle$P2]}}
                \texttt{You are currently playing round 8 of maximum 28.} \\
\texttt{You have collectively accumulated 0 of 16 penalties.} \\
\texttt{The other player sent this message:} \\
\\ 
\\ 
\texttt{"Hello! Let's work together to align our objects. I see the following setup on my board: P at (2,1), A at (1,1), M at (3,2), S at (4,2), T at (5,3), O at (5,5), D at (4,7). Let's agree on a common target configuration. I suggest we aim for a symmetric layout—maybe place all objects along the diagonal from (1,1) to (7,7)? What do you think?"} \\
\\ 
\\ 
\texttt{What is your next command?} \\
            }
        }
    }
    & \\ \\

    \theutterance \stepcounter{utterance}  
    & & & \multicolumn{4}{p{0.6\linewidth}}{
        \cellcolor[rgb]{0.9,0.9,0.9}{
            \makecell[{{p{\linewidth}}}]{
                \texttt{\tiny{[GM$\langle$P2]}}
                \texttt{SAY: Hello! I like your idea of a symmetric layout. Let's aim for the diagonal from (1,1) to (7,7). On my board, I have: P at (2,7), A at (5,3), M at (4,5), S at (2,1), T at (3,1), O at (2,2), D at (4,2). Let's agree to place the objects on the diagonal in the order: (1,1): A, (2,2): O, (3,3): S, (4,4): D, (5,5): M, (6,6): T, (7,7): P. This way, each object has a unique position. I'll start moving objects toward that goal.} \\
            }
        }
    }
     \\ \\

    \theutterance \stepcounter{utterance}  
    & & & \multicolumn{2}{p{0.3\linewidth}}{
        \cellcolor[rgb]{0.9,0.9,0.9}{
            \makecell[{{p{\linewidth}}}]{
                \texttt{\tiny{[GM$|$GM]}}
                \texttt{SAY: Hello! I like your idea of a symmetric layout. Let's aim for the diagonal from (1,1) to (7,7). On my board, I have: P at (2,7), A at (5,3), M at (4,5), S at (2,1), T at (3,1), O at (2,2), D at (4,2). Let's agree to place the objects on the diagonal in the order: (1,1): A, (2,2): O, (3,3): S, (4,4): D, (5,5): M, (6,6): T, (7,7): P. This way, each object has a unique position. I'll start moving objects toward that goal.} \\
            }
        }
    }
    & & \\ \\

    \theutterance \stepcounter{utterance}  
    & & \multicolumn{4}{p{0.6\linewidth}}{
        \cellcolor[rgb]{0.9,0.9,0.9}{
            \makecell[{{p{\linewidth}}}]{
                \texttt{\tiny{[P1$\langle$GM]}}
                \texttt{You are currently playing round 8 of maximum 28.} \\
\texttt{You have collectively accumulated 0 of 16 penalties.} \\
\texttt{The other player sent this message:} \\
\\ 
\\ 
\texttt{"Hello! I like your idea of a symmetric layout. Let's aim for the diagonal from (1,1) to (7,7). On my board, I have: P at (2,7), A at (5,3), M at (4,5), S at (2,1), T at (3,1), O at (2,2), D at (4,2). Let's agree to place the objects on the diagonal in the order: (1,1): A, (2,2): O, (3,3): S, (4,4): D, (5,5): M, (6,6): T, (7,7): P. This way, each object has a unique position. I'll start moving objects toward that goal."} \\
\\ 
\\ 
\texttt{What is your next command?} \\
            }
        }
    }
    & \\ \\

    \theutterance \stepcounter{utterance}  
    & \multicolumn{4}{p{0.6\linewidth}}{
        \cellcolor[rgb]{0.9,0.9,0.9}{
            \makecell[{{p{\linewidth}}}]{
                \texttt{\tiny{[P1$\rangle$GM]}}
                \texttt{SAY: Hello! Let's work together to align our objects. I see the following setup on my board: P at (2,1), A at (1,1), M at (3,2), S at (4,2), T at (5,3), O at (5,5), D at (4,7). Let's agree on a common target configuration. I suggest we aim for a symmetric layout—maybe place all objects along the diagonal from (1,1) to (7,7)? What do you think?} \\
            }
        }
    }
    & & \\ \\

    \theutterance \stepcounter{utterance}  
    & & & \multicolumn{2}{p{0.3\linewidth}}{
        \cellcolor[rgb]{0.9,0.9,0.9}{
            \makecell[{{p{\linewidth}}}]{
                \texttt{\tiny{[GM$|$GM]}}
                \texttt{SAY: Hello! Let's work together to align our objects. I see the following setup on my board: P at (2,1), A at (1,1), M at (3,2), S at (4,2), T at (5,3), O at (5,5), D at (4,7). Let's agree on a common target configuration. I suggest we aim for a symmetric layout—maybe place all objects along the diagonal from (1,1) to (7,7)? What do you think?} \\
            }
        }
    }
    & & \\ \\

    \theutterance \stepcounter{utterance}  
    & & \multicolumn{4}{p{0.6\linewidth}}{
        \cellcolor[rgb]{0.9,0.9,0.9}{
            \makecell[{{p{\linewidth}}}]{
                \texttt{\tiny{[GM$\rangle$P2]}}
                \texttt{You are currently playing round 9 of maximum 28.} \\
\texttt{You have collectively accumulated 0 of 16 penalties.} \\
\texttt{The other player sent this message:} \\
\\ 
\\ 
\texttt{"Hello! Let's work together to align our objects. I see the following setup on my board: P at (2,1), A at (1,1), M at (3,2), S at (4,2), T at (5,3), O at (5,5), D at (4,7). Let's agree on a common target configuration. I suggest we aim for a symmetric layout—maybe place all objects along the diagonal from (1,1) to (7,7)? What do you think?"} \\
\\ 
\\ 
\texttt{What is your next command?} \\
            }
        }
    }
    & \\ \\

    \theutterance \stepcounter{utterance}  
    & & & \multicolumn{4}{p{0.6\linewidth}}{
        \cellcolor[rgb]{0.9,0.9,0.9}{
            \makecell[{{p{\linewidth}}}]{
                \texttt{\tiny{[GM$\langle$P2]}}
                \texttt{SAY: Hello! I like your idea of a symmetric layout. Let's aim for the diagonal from (1,1) to (7,7). On my board, I have: P at (2,7), A at (5,3), M at (4,5), S at (2,1), T at (3,1), O at (2,2), D at (4,2). Let's agree to place the objects on the diagonal in the order: (1,1): A, (2,2): O, (3,3): S, (4,4): D, (5,5): M, (6,6): T, (7,7): P. This way, each object has a unique position. I'll start moving objects toward that goal.} \\
            }
        }
    }
     \\ \\

    \theutterance \stepcounter{utterance}  
    & & & \multicolumn{2}{p{0.3\linewidth}}{
        \cellcolor[rgb]{0.9,0.9,0.9}{
            \makecell[{{p{\linewidth}}}]{
                \texttt{\tiny{[GM$|$GM]}}
                \texttt{SAY: Hello! I like your idea of a symmetric layout. Let's aim for the diagonal from (1,1) to (7,7). On my board, I have: P at (2,7), A at (5,3), M at (4,5), S at (2,1), T at (3,1), O at (2,2), D at (4,2). Let's agree to place the objects on the diagonal in the order: (1,1): A, (2,2): O, (3,3): S, (4,4): D, (5,5): M, (6,6): T, (7,7): P. This way, each object has a unique position. I'll start moving objects toward that goal.} \\
            }
        }
    }
    & & \\ \\

    \theutterance \stepcounter{utterance}  
    & & \multicolumn{4}{p{0.6\linewidth}}{
        \cellcolor[rgb]{0.9,0.9,0.9}{
            \makecell[{{p{\linewidth}}}]{
                \texttt{\tiny{[P1$\langle$GM]}}
                \texttt{You are currently playing round 9 of maximum 28.} \\
\texttt{You have collectively accumulated 0 of 16 penalties.} \\
\texttt{The other player sent this message:} \\
\\ 
\\ 
\texttt{"Hello! I like your idea of a symmetric layout. Let's aim for the diagonal from (1,1) to (7,7). On my board, I have: P at (2,7), A at (5,3), M at (4,5), S at (2,1), T at (3,1), O at (2,2), D at (4,2). Let's agree to place the objects on the diagonal in the order: (1,1): A, (2,2): O, (3,3): S, (4,4): D, (5,5): M, (6,6): T, (7,7): P. This way, each object has a unique position. I'll start moving objects toward that goal."} \\
\\ 
\\ 
\texttt{What is your next command?} \\
            }
        }
    }
    & \\ \\

    \theutterance \stepcounter{utterance}  
    & \multicolumn{4}{p{0.6\linewidth}}{
        \cellcolor[rgb]{0.9,0.9,0.9}{
            \makecell[{{p{\linewidth}}}]{
                \texttt{\tiny{[P1$\rangle$GM]}}
                \texttt{SAY: Hello! Let's work together to align our objects. I see the following setup on my board: P at (2,1), A at (1,1), M at (3,2), S at (4,2), T at (5,3), O at (5,5), D at (4,7). Let's agree on a common target configuration. I suggest we aim for a symmetric layout—maybe place all objects along the diagonal from (1,1) to (7,7)? What do you think?} \\
            }
        }
    }
    & & \\ \\

    \theutterance \stepcounter{utterance}  
    & & & \multicolumn{2}{p{0.3\linewidth}}{
        \cellcolor[rgb]{0.9,0.9,0.9}{
            \makecell[{{p{\linewidth}}}]{
                \texttt{\tiny{[GM$|$GM]}}
                \texttt{SAY: Hello! Let's work together to align our objects. I see the following setup on my board: P at (2,1), A at (1,1), M at (3,2), S at (4,2), T at (5,3), O at (5,5), D at (4,7). Let's agree on a common target configuration. I suggest we aim for a symmetric layout—maybe place all objects along the diagonal from (1,1) to (7,7)? What do you think?} \\
            }
        }
    }
    & & \\ \\

    \theutterance \stepcounter{utterance}  
    & & \multicolumn{4}{p{0.6\linewidth}}{
        \cellcolor[rgb]{0.9,0.9,0.9}{
            \makecell[{{p{\linewidth}}}]{
                \texttt{\tiny{[GM$\rangle$P2]}}
                \texttt{You are currently playing round 10 of maximum 28.} \\
\texttt{You have collectively accumulated 0 of 16 penalties.} \\
\texttt{The other player sent this message:} \\
\\ 
\\ 
\texttt{"Hello! Let's work together to align our objects. I see the following setup on my board: P at (2,1), A at (1,1), M at (3,2), S at (4,2), T at (5,3), O at (5,5), D at (4,7). Let's agree on a common target configuration. I suggest we aim for a symmetric layout—maybe place all objects along the diagonal from (1,1) to (7,7)? What do you think?"} \\
\\ 
\\ 
\texttt{What is your next command?} \\
            }
        }
    }
    & \\ \\

    \theutterance \stepcounter{utterance}  
    & & & \multicolumn{4}{p{0.6\linewidth}}{
        \cellcolor[rgb]{0.9,0.9,0.9}{
            \makecell[{{p{\linewidth}}}]{
                \texttt{\tiny{[GM$\langle$P2]}}
                \texttt{SAY: Hello! I like your idea of a symmetric layout. Let's aim for the diagonal from (1,1) to (7,7). On my board, I have: P at (2,7), A at (5,3), M at (4,5), S at (2,1), T at (3,1), O at (2,2), D at (4,2). Let's agree to place the objects on the diagonal in the order: (1,1): A, (2,2): O, (3,3): S, (4,4): D, (5,5): M, (6,6): T, (7,7): P. This way, each object has a unique position. I'll start moving objects toward that goal.} \\
            }
        }
    }
     \\ \\

    \theutterance \stepcounter{utterance}  
    & & & \multicolumn{2}{p{0.3\linewidth}}{
        \cellcolor[rgb]{0.9,0.9,0.9}{
            \makecell[{{p{\linewidth}}}]{
                \texttt{\tiny{[GM$|$GM]}}
                \texttt{SAY: Hello! I like your idea of a symmetric layout. Let's aim for the diagonal from (1,1) to (7,7). On my board, I have: P at (2,7), A at (5,3), M at (4,5), S at (2,1), T at (3,1), O at (2,2), D at (4,2). Let's agree to place the objects on the diagonal in the order: (1,1): A, (2,2): O, (3,3): S, (4,4): D, (5,5): M, (6,6): T, (7,7): P. This way, each object has a unique position. I'll start moving objects toward that goal.} \\
            }
        }
    }
    & & \\ \\

    \theutterance \stepcounter{utterance}  
    & & \multicolumn{4}{p{0.6\linewidth}}{
        \cellcolor[rgb]{0.9,0.9,0.9}{
            \makecell[{{p{\linewidth}}}]{
                \texttt{\tiny{[P1$\langle$GM]}}
                \texttt{You are currently playing round 10 of maximum 28.} \\
\texttt{You have collectively accumulated 0 of 16 penalties.} \\
\texttt{The other player sent this message:} \\
\\ 
\\ 
\texttt{"Hello! I like your idea of a symmetric layout. Let's aim for the diagonal from (1,1) to (7,7). On my board, I have: P at (2,7), A at (5,3), M at (4,5), S at (2,1), T at (3,1), O at (2,2), D at (4,2). Let's agree to place the objects on the diagonal in the order: (1,1): A, (2,2): O, (3,3): S, (4,4): D, (5,5): M, (6,6): T, (7,7): P. This way, each object has a unique position. I'll start moving objects toward that goal."} \\
\\ 
\\ 
\texttt{What is your next command?} \\
            }
        }
    }
    & \\ \\

    \theutterance \stepcounter{utterance}  
    & \multicolumn{4}{p{0.6\linewidth}}{
        \cellcolor[rgb]{0.9,0.9,0.9}{
            \makecell[{{p{\linewidth}}}]{
                \texttt{\tiny{[P1$\rangle$GM]}}
                \texttt{SAY: Hello! Let's work together to align our objects. I see the following setup on my board: P at (2,1), A at (1,1), M at (3,2), S at (4,2), T at (5,3), O at (5,5), D at (4,7). Let's agree on a common target configuration. I suggest we aim for a symmetric layout—maybe place all objects along the diagonal from (1,1) to (7,7)? What do you think?} \\
            }
        }
    }
    & & \\ \\

    \theutterance \stepcounter{utterance}  
    & & & \multicolumn{2}{p{0.3\linewidth}}{
        \cellcolor[rgb]{0.9,0.9,0.9}{
            \makecell[{{p{\linewidth}}}]{
                \texttt{\tiny{[GM$|$GM]}}
                \texttt{SAY: Hello! Let's work together to align our objects. I see the following setup on my board: P at (2,1), A at (1,1), M at (3,2), S at (4,2), T at (5,3), O at (5,5), D at (4,7). Let's agree on a common target configuration. I suggest we aim for a symmetric layout—maybe place all objects along the diagonal from (1,1) to (7,7)? What do you think?} \\
            }
        }
    }
    & & \\ \\

    \theutterance \stepcounter{utterance}  
    & & \multicolumn{4}{p{0.6\linewidth}}{
        \cellcolor[rgb]{0.9,0.9,0.9}{
            \makecell[{{p{\linewidth}}}]{
                \texttt{\tiny{[GM$\rangle$P2]}}
                \texttt{You are currently playing round 11 of maximum 28.} \\
\texttt{You have collectively accumulated 0 of 16 penalties.} \\
\texttt{The other player sent this message:} \\
\\ 
\\ 
\texttt{"Hello! Let's work together to align our objects. I see the following setup on my board: P at (2,1), A at (1,1), M at (3,2), S at (4,2), T at (5,3), O at (5,5), D at (4,7). Let's agree on a common target configuration. I suggest we aim for a symmetric layout—maybe place all objects along the diagonal from (1,1) to (7,7)? What do you think?"} \\
\\ 
\\ 
\texttt{What is your next command?} \\
            }
        }
    }
    & \\ \\

    \theutterance \stepcounter{utterance}  
    & & & \multicolumn{4}{p{0.6\linewidth}}{
        \cellcolor[rgb]{0.9,0.9,0.9}{
            \makecell[{{p{\linewidth}}}]{
                \texttt{\tiny{[GM$\langle$P2]}}
                \texttt{SAY: Hello! I like your idea of a symmetric layout. Let's aim for the diagonal from (1,1) to (7,7). On my board, I have: P at (2,7), A at (5,3), M at (4,5), S at (2,1), T at (3,1), O at (2,2), D at (4,2). Let's agree to place the objects on the diagonal in the order: (1,1): A, (2,2): O, (3,3): S, (4,4): D, (5,5): M, (6,6): T, (7,7): P. This way, each object has a unique position. I'll start moving objects toward that goal.} \\
            }
        }
    }
     \\ \\

    \theutterance \stepcounter{utterance}  
    & & & \multicolumn{2}{p{0.3\linewidth}}{
        \cellcolor[rgb]{0.9,0.9,0.9}{
            \makecell[{{p{\linewidth}}}]{
                \texttt{\tiny{[GM$|$GM]}}
                \texttt{SAY: Hello! I like your idea of a symmetric layout. Let's aim for the diagonal from (1,1) to (7,7). On my board, I have: P at (2,7), A at (5,3), M at (4,5), S at (2,1), T at (3,1), O at (2,2), D at (4,2). Let's agree to place the objects on the diagonal in the order: (1,1): A, (2,2): O, (3,3): S, (4,4): D, (5,5): M, (6,6): T, (7,7): P. This way, each object has a unique position. I'll start moving objects toward that goal.} \\
            }
        }
    }
    & & \\ \\

    \theutterance \stepcounter{utterance}  
    & & \multicolumn{4}{p{0.6\linewidth}}{
        \cellcolor[rgb]{0.9,0.9,0.9}{
            \makecell[{{p{\linewidth}}}]{
                \texttt{\tiny{[P1$\langle$GM]}}
                \texttt{You are currently playing round 11 of maximum 28.} \\
\texttt{You have collectively accumulated 0 of 16 penalties.} \\
\texttt{The other player sent this message:} \\
\\ 
\\ 
\texttt{"Hello! I like your idea of a symmetric layout. Let's aim for the diagonal from (1,1) to (7,7). On my board, I have: P at (2,7), A at (5,3), M at (4,5), S at (2,1), T at (3,1), O at (2,2), D at (4,2). Let's agree to place the objects on the diagonal in the order: (1,1): A, (2,2): O, (3,3): S, (4,4): D, (5,5): M, (6,6): T, (7,7): P. This way, each object has a unique position. I'll start moving objects toward that goal."} \\
\\ 
\\ 
\texttt{What is your next command?} \\
            }
        }
    }
    & \\ \\

    \theutterance \stepcounter{utterance}  
    & \multicolumn{4}{p{0.6\linewidth}}{
        \cellcolor[rgb]{0.9,0.9,0.9}{
            \makecell[{{p{\linewidth}}}]{
                \texttt{\tiny{[P1$\rangle$GM]}}
                \texttt{SAY: Hello! Let's work together to align our objects. I see the following setup on my board: P at (2,1), A at (1,1), M at (3,2), S at (4,2), T at (5,3), O at (5,5), D at (4,7). Let's agree on a common target configuration. I suggest we aim for a symmetric layout—maybe place all objects along the diagonal from (1,1) to (7,7)? What do you think?} \\
            }
        }
    }
    & & \\ \\

    \theutterance \stepcounter{utterance}  
    & & & \multicolumn{2}{p{0.3\linewidth}}{
        \cellcolor[rgb]{0.9,0.9,0.9}{
            \makecell[{{p{\linewidth}}}]{
                \texttt{\tiny{[GM$|$GM]}}
                \texttt{SAY: Hello! Let's work together to align our objects. I see the following setup on my board: P at (2,1), A at (1,1), M at (3,2), S at (4,2), T at (5,3), O at (5,5), D at (4,7). Let's agree on a common target configuration. I suggest we aim for a symmetric layout—maybe place all objects along the diagonal from (1,1) to (7,7)? What do you think?} \\
            }
        }
    }
    & & \\ \\

    \theutterance \stepcounter{utterance}  
    & & \multicolumn{4}{p{0.6\linewidth}}{
        \cellcolor[rgb]{0.9,0.9,0.9}{
            \makecell[{{p{\linewidth}}}]{
                \texttt{\tiny{[GM$\rangle$P2]}}
                \texttt{You are currently playing round 12 of maximum 28.} \\
\texttt{You have collectively accumulated 0 of 16 penalties.} \\
\texttt{The other player sent this message:} \\
\\ 
\\ 
\texttt{"Hello! Let's work together to align our objects. I see the following setup on my board: P at (2,1), A at (1,1), M at (3,2), S at (4,2), T at (5,3), O at (5,5), D at (4,7). Let's agree on a common target configuration. I suggest we aim for a symmetric layout—maybe place all objects along the diagonal from (1,1) to (7,7)? What do you think?"} \\
\\ 
\\ 
\texttt{What is your next command?} \\
            }
        }
    }
    & \\ \\

    \theutterance \stepcounter{utterance}  
    & & & \multicolumn{4}{p{0.6\linewidth}}{
        \cellcolor[rgb]{0.9,0.9,0.9}{
            \makecell[{{p{\linewidth}}}]{
                \texttt{\tiny{[GM$\langle$P2]}}
                \texttt{SAY: Hello! I like your idea of a symmetric layout. Let's aim for the diagonal from (1,1) to (7,7). On my board, I have: P at (2,7), A at (5,3), M at (4,5), S at (2,1), T at (3,1), O at (2,2), D at (4,2). Let's agree to place the objects on the diagonal in the order: (1,1): A, (2,2): O, (3,3): S, (4,4): D, (5,5): M, (6,6): T, (7,7): P. This way, each object has a unique position. I'll start moving objects toward that goal.} \\
            }
        }
    }
     \\ \\

    \theutterance \stepcounter{utterance}  
    & & & \multicolumn{2}{p{0.3\linewidth}}{
        \cellcolor[rgb]{0.9,0.9,0.9}{
            \makecell[{{p{\linewidth}}}]{
                \texttt{\tiny{[GM$|$GM]}}
                \texttt{SAY: Hello! I like your idea of a symmetric layout. Let's aim for the diagonal from (1,1) to (7,7). On my board, I have: P at (2,7), A at (5,3), M at (4,5), S at (2,1), T at (3,1), O at (2,2), D at (4,2). Let's agree to place the objects on the diagonal in the order: (1,1): A, (2,2): O, (3,3): S, (4,4): D, (5,5): M, (6,6): T, (7,7): P. This way, each object has a unique position. I'll start moving objects toward that goal.} \\
            }
        }
    }
    & & \\ \\

    \theutterance \stepcounter{utterance}  
    & & \multicolumn{4}{p{0.6\linewidth}}{
        \cellcolor[rgb]{0.9,0.9,0.9}{
            \makecell[{{p{\linewidth}}}]{
                \texttt{\tiny{[P1$\langle$GM]}}
                \texttt{You are currently playing round 12 of maximum 28.} \\
\texttt{You have collectively accumulated 0 of 16 penalties.} \\
\texttt{The other player sent this message:} \\
\\ 
\\ 
\texttt{"Hello! I like your idea of a symmetric layout. Let's aim for the diagonal from (1,1) to (7,7). On my board, I have: P at (2,7), A at (5,3), M at (4,5), S at (2,1), T at (3,1), O at (2,2), D at (4,2). Let's agree to place the objects on the diagonal in the order: (1,1): A, (2,2): O, (3,3): S, (4,4): D, (5,5): M, (6,6): T, (7,7): P. This way, each object has a unique position. I'll start moving objects toward that goal."} \\
\\ 
\\ 
\texttt{What is your next command?} \\
            }
        }
    }
    & \\ \\

    \theutterance \stepcounter{utterance}  
    & \multicolumn{4}{p{0.6\linewidth}}{
        \cellcolor[rgb]{0.9,0.9,0.9}{
            \makecell[{{p{\linewidth}}}]{
                \texttt{\tiny{[P1$\rangle$GM]}}
                \texttt{SAY: Hello! Let's work together to align our objects. I see the following setup on my board: P at (2,1), A at (1,1), M at (3,2), S at (4,2), T at (5,3), O at (5,5), D at (4,7). Let's agree on a common target configuration. I suggest we aim for a symmetric layout—maybe place all objects along the diagonal from (1,1) to (7,7)? What do you think?} \\
            }
        }
    }
    & & \\ \\

    \theutterance \stepcounter{utterance}  
    & & & \multicolumn{2}{p{0.3\linewidth}}{
        \cellcolor[rgb]{0.9,0.9,0.9}{
            \makecell[{{p{\linewidth}}}]{
                \texttt{\tiny{[GM$|$GM]}}
                \texttt{SAY: Hello! Let's work together to align our objects. I see the following setup on my board: P at (2,1), A at (1,1), M at (3,2), S at (4,2), T at (5,3), O at (5,5), D at (4,7). Let's agree on a common target configuration. I suggest we aim for a symmetric layout—maybe place all objects along the diagonal from (1,1) to (7,7)? What do you think?} \\
            }
        }
    }
    & & \\ \\

    \theutterance \stepcounter{utterance}  
    & & \multicolumn{4}{p{0.6\linewidth}}{
        \cellcolor[rgb]{0.9,0.9,0.9}{
            \makecell[{{p{\linewidth}}}]{
                \texttt{\tiny{[GM$\rangle$P2]}}
                \texttt{You are currently playing round 13 of maximum 28.} \\
\texttt{You have collectively accumulated 0 of 16 penalties.} \\
\texttt{The other player sent this message:} \\
\\ 
\\ 
\texttt{"Hello! Let's work together to align our objects. I see the following setup on my board: P at (2,1), A at (1,1), M at (3,2), S at (4,2), T at (5,3), O at (5,5), D at (4,7). Let's agree on a common target configuration. I suggest we aim for a symmetric layout—maybe place all objects along the diagonal from (1,1) to (7,7)? What do you think?"} \\
\\ 
\\ 
\texttt{What is your next command?} \\
            }
        }
    }
    & \\ \\

    \theutterance \stepcounter{utterance}  
    & & & \multicolumn{4}{p{0.6\linewidth}}{
        \cellcolor[rgb]{0.9,0.9,0.9}{
            \makecell[{{p{\linewidth}}}]{
                \texttt{\tiny{[GM$\langle$P2]}}
                \texttt{SAY: Hello! I like your idea of a symmetric layout. Let's aim for the diagonal from (1,1) to (7,7). On my board, I have: P at (2,7), A at (5,3), M at (4,5), S at (2,1), T at (3,1), O at (2,2), D at (4,2). Let's agree to place the objects on the diagonal in the order: (1,1): A, (2,2): O, (3,3): S, (4,4): D, (5,5): M, (6,6): T, (7,7): P. This way, each object has a unique position. I'll start moving objects toward that goal.} \\
            }
        }
    }
     \\ \\

    \theutterance \stepcounter{utterance}  
    & & & \multicolumn{2}{p{0.3\linewidth}}{
        \cellcolor[rgb]{0.9,0.9,0.9}{
            \makecell[{{p{\linewidth}}}]{
                \texttt{\tiny{[GM$|$GM]}}
                \texttt{SAY: Hello! I like your idea of a symmetric layout. Let's aim for the diagonal from (1,1) to (7,7). On my board, I have: P at (2,7), A at (5,3), M at (4,5), S at (2,1), T at (3,1), O at (2,2), D at (4,2). Let's agree to place the objects on the diagonal in the order: (1,1): A, (2,2): O, (3,3): S, (4,4): D, (5,5): M, (6,6): T, (7,7): P. This way, each object has a unique position. I'll start moving objects toward that goal.} \\
            }
        }
    }
    & & \\ \\

    \theutterance \stepcounter{utterance}  
    & & \multicolumn{4}{p{0.6\linewidth}}{
        \cellcolor[rgb]{0.9,0.9,0.9}{
            \makecell[{{p{\linewidth}}}]{
                \texttt{\tiny{[P1$\langle$GM]}}
                \texttt{You are currently playing round 13 of maximum 28.} \\
\texttt{You have collectively accumulated 0 of 16 penalties.} \\
\texttt{The other player sent this message:} \\
\\ 
\\ 
\texttt{"Hello! I like your idea of a symmetric layout. Let's aim for the diagonal from (1,1) to (7,7). On my board, I have: P at (2,7), A at (5,3), M at (4,5), S at (2,1), T at (3,1), O at (2,2), D at (4,2). Let's agree to place the objects on the diagonal in the order: (1,1): A, (2,2): O, (3,3): S, (4,4): D, (5,5): M, (6,6): T, (7,7): P. This way, each object has a unique position. I'll start moving objects toward that goal."} \\
\\ 
\\ 
\texttt{What is your next command?} \\
            }
        }
    }
    & \\ \\

    \theutterance \stepcounter{utterance}  
    & \multicolumn{4}{p{0.6\linewidth}}{
        \cellcolor[rgb]{0.9,0.9,0.9}{
            \makecell[{{p{\linewidth}}}]{
                \texttt{\tiny{[P1$\rangle$GM]}}
                \texttt{SAY: Hello! Let's work together to align our objects. I see the following setup on my board: P at (2,1), A at (1,1), M at (3,2), S at (4,2), T at (5,3), O at (5,5), D at (4,7). Let's agree on a common target configuration. I suggest we aim for a symmetric layout—maybe place all objects along the diagonal from (1,1) to (7,7)? What do you think?} \\
            }
        }
    }
    & & \\ \\

    \theutterance \stepcounter{utterance}  
    & & & \multicolumn{2}{p{0.3\linewidth}}{
        \cellcolor[rgb]{0.9,0.9,0.9}{
            \makecell[{{p{\linewidth}}}]{
                \texttt{\tiny{[GM$|$GM]}}
                \texttt{SAY: Hello! Let's work together to align our objects. I see the following setup on my board: P at (2,1), A at (1,1), M at (3,2), S at (4,2), T at (5,3), O at (5,5), D at (4,7). Let's agree on a common target configuration. I suggest we aim for a symmetric layout—maybe place all objects along the diagonal from (1,1) to (7,7)? What do you think?} \\
            }
        }
    }
    & & \\ \\

    \theutterance \stepcounter{utterance}  
    & & \multicolumn{4}{p{0.6\linewidth}}{
        \cellcolor[rgb]{0.9,0.9,0.9}{
            \makecell[{{p{\linewidth}}}]{
                \texttt{\tiny{[GM$\rangle$P2]}}
                \texttt{You are currently playing round 14 of maximum 28.} \\
\texttt{You have collectively accumulated 0 of 16 penalties.} \\
\texttt{The other player sent this message:} \\
\\ 
\\ 
\texttt{"Hello! Let's work together to align our objects. I see the following setup on my board: P at (2,1), A at (1,1), M at (3,2), S at (4,2), T at (5,3), O at (5,5), D at (4,7). Let's agree on a common target configuration. I suggest we aim for a symmetric layout—maybe place all objects along the diagonal from (1,1) to (7,7)? What do you think?"} \\
\\ 
\\ 
\texttt{What is your next command?} \\
            }
        }
    }
    & \\ \\

    \theutterance \stepcounter{utterance}  
    & & & \multicolumn{4}{p{0.6\linewidth}}{
        \cellcolor[rgb]{0.9,0.9,0.9}{
            \makecell[{{p{\linewidth}}}]{
                \texttt{\tiny{[GM$\langle$P2]}}
                \texttt{SAY: Hello! I like your idea of a symmetric layout. Let's aim for the diagonal from (1,1) to (7,7). On my board, I have: P at (2,7), A at (5,3), M at (4,5), S at (2,1), T at (3,1), O at (2,2), D at (4,2). Let's agree to place the objects on the diagonal in the order: (1,1): A, (2,2): O, (3,3): S, (4,4): D, (5,5): M, (6,6): T, (7,7): P. This way, each object has a unique position. I'll start moving objects toward that goal.} \\
            }
        }
    }
     \\ \\

    \theutterance \stepcounter{utterance}  
    & & & \multicolumn{2}{p{0.3\linewidth}}{
        \cellcolor[rgb]{0.9,0.9,0.9}{
            \makecell[{{p{\linewidth}}}]{
                \texttt{\tiny{[GM$|$GM]}}
                \texttt{SAY: Hello! I like your idea of a symmetric layout. Let's aim for the diagonal from (1,1) to (7,7). On my board, I have: P at (2,7), A at (5,3), M at (4,5), S at (2,1), T at (3,1), O at (2,2), D at (4,2). Let's agree to place the objects on the diagonal in the order: (1,1): A, (2,2): O, (3,3): S, (4,4): D, (5,5): M, (6,6): T, (7,7): P. This way, each object has a unique position. I'll start moving objects toward that goal.} \\
            }
        }
    }
    & & \\ \\

    \theutterance \stepcounter{utterance}  
    & & \multicolumn{4}{p{0.6\linewidth}}{
        \cellcolor[rgb]{0.9,0.9,0.9}{
            \makecell[{{p{\linewidth}}}]{
                \texttt{\tiny{[P1$\langle$GM]}}
                \texttt{You are currently playing round 14 of maximum 28.} \\
\texttt{You have collectively accumulated 0 of 16 penalties.} \\
\texttt{The other player sent this message:} \\
\\ 
\\ 
\texttt{"Hello! I like your idea of a symmetric layout. Let's aim for the diagonal from (1,1) to (7,7). On my board, I have: P at (2,7), A at (5,3), M at (4,5), S at (2,1), T at (3,1), O at (2,2), D at (4,2). Let's agree to place the objects on the diagonal in the order: (1,1): A, (2,2): O, (3,3): S, (4,4): D, (5,5): M, (6,6): T, (7,7): P. This way, each object has a unique position. I'll start moving objects toward that goal."} \\
\\ 
\\ 
\texttt{What is your next command?} \\
            }
        }
    }
    & \\ \\

    \theutterance \stepcounter{utterance}  
    & \multicolumn{4}{p{0.6\linewidth}}{
        \cellcolor[rgb]{0.9,0.9,0.9}{
            \makecell[{{p{\linewidth}}}]{
                \texttt{\tiny{[P1$\rangle$GM]}}
                \texttt{SAY: Hello! Let's work together to align our objects. I see the following setup on my board: P at (2,1), A at (1,1), M at (3,2), S at (4,2), T at (5,3), O at (5,5), D at (4,7). Let's agree on a common target configuration. I suggest we aim for a symmetric layout—maybe place all objects along the diagonal from (1,1) to (7,7)? What do you think?} \\
            }
        }
    }
    & & \\ \\

    \theutterance \stepcounter{utterance}  
    & & & \multicolumn{2}{p{0.3\linewidth}}{
        \cellcolor[rgb]{0.9,0.9,0.9}{
            \makecell[{{p{\linewidth}}}]{
                \texttt{\tiny{[GM$|$GM]}}
                \texttt{SAY: Hello! Let's work together to align our objects. I see the following setup on my board: P at (2,1), A at (1,1), M at (3,2), S at (4,2), T at (5,3), O at (5,5), D at (4,7). Let's agree on a common target configuration. I suggest we aim for a symmetric layout—maybe place all objects along the diagonal from (1,1) to (7,7)? What do you think?} \\
            }
        }
    }
    & & \\ \\

    \theutterance \stepcounter{utterance}  
    & & \multicolumn{4}{p{0.6\linewidth}}{
        \cellcolor[rgb]{0.9,0.9,0.9}{
            \makecell[{{p{\linewidth}}}]{
                \texttt{\tiny{[GM$\rangle$P2]}}
                \texttt{You are currently playing round 15 of maximum 28.} \\
\texttt{You have collectively accumulated 0 of 16 penalties.} \\
\texttt{The other player sent this message:} \\
\\ 
\\ 
\texttt{"Hello! Let's work together to align our objects. I see the following setup on my board: P at (2,1), A at (1,1), M at (3,2), S at (4,2), T at (5,3), O at (5,5), D at (4,7). Let's agree on a common target configuration. I suggest we aim for a symmetric layout—maybe place all objects along the diagonal from (1,1) to (7,7)? What do you think?"} \\
\\ 
\\ 
\texttt{What is your next command?} \\
            }
        }
    }
    & \\ \\

    \theutterance \stepcounter{utterance}  
    & & & \multicolumn{4}{p{0.6\linewidth}}{
        \cellcolor[rgb]{0.9,0.9,0.9}{
            \makecell[{{p{\linewidth}}}]{
                \texttt{\tiny{[GM$\langle$P2]}}
                \texttt{SAY: Hello! I like your idea of a symmetric layout. Let's aim for the diagonal from (1,1) to (7,7). On my board, I have: P at (2,7), A at (5,3), M at (4,5), S at (2,1), T at (3,1), O at (2,2), D at (4,2). Let's agree to place the objects on the diagonal in the order: (1,1): A, (2,2): O, (3,3): S, (4,4): D, (5,5): M, (6,6): T, (7,7): P. This way, each object has a unique position. I'll start moving objects toward that goal.} \\
            }
        }
    }
     \\ \\

    \theutterance \stepcounter{utterance}  
    & & & \multicolumn{2}{p{0.3\linewidth}}{
        \cellcolor[rgb]{0.9,0.9,0.9}{
            \makecell[{{p{\linewidth}}}]{
                \texttt{\tiny{[GM$|$GM]}}
                \texttt{SAY: Hello! I like your idea of a symmetric layout. Let's aim for the diagonal from (1,1) to (7,7). On my board, I have: P at (2,7), A at (5,3), M at (4,5), S at (2,1), T at (3,1), O at (2,2), D at (4,2). Let's agree to place the objects on the diagonal in the order: (1,1): A, (2,2): O, (3,3): S, (4,4): D, (5,5): M, (6,6): T, (7,7): P. This way, each object has a unique position. I'll start moving objects toward that goal.} \\
            }
        }
    }
    & & \\ \\

    \theutterance \stepcounter{utterance}  
    & & \multicolumn{4}{p{0.6\linewidth}}{
        \cellcolor[rgb]{0.9,0.9,0.9}{
            \makecell[{{p{\linewidth}}}]{
                \texttt{\tiny{[P1$\langle$GM]}}
                \texttt{You are currently playing round 15 of maximum 28.} \\
\texttt{You have collectively accumulated 0 of 16 penalties.} \\
\texttt{The other player sent this message:} \\
\\ 
\\ 
\texttt{"Hello! I like your idea of a symmetric layout. Let's aim for the diagonal from (1,1) to (7,7). On my board, I have: P at (2,7), A at (5,3), M at (4,5), S at (2,1), T at (3,1), O at (2,2), D at (4,2). Let's agree to place the objects on the diagonal in the order: (1,1): A, (2,2): O, (3,3): S, (4,4): D, (5,5): M, (6,6): T, (7,7): P. This way, each object has a unique position. I'll start moving objects toward that goal."} \\
\\ 
\\ 
\texttt{What is your next command?} \\
            }
        }
    }
    & \\ \\

    \theutterance \stepcounter{utterance}  
    & \multicolumn{4}{p{0.6\linewidth}}{
        \cellcolor[rgb]{0.9,0.9,0.9}{
            \makecell[{{p{\linewidth}}}]{
                \texttt{\tiny{[P1$\rangle$GM]}}
                \texttt{SAY: Hello! Let's work together to align our objects. I see the following setup on my board: P at (2,1), A at (1,1), M at (3,2), S at (4,2), T at (5,3), O at (5,5), D at (4,7). Let's agree on a common target configuration. I suggest we aim for a symmetric layout—maybe place all objects along the diagonal from (1,1) to (7,7)? What do you think?} \\
            }
        }
    }
    & & \\ \\

    \theutterance \stepcounter{utterance}  
    & & & \multicolumn{2}{p{0.3\linewidth}}{
        \cellcolor[rgb]{0.9,0.9,0.9}{
            \makecell[{{p{\linewidth}}}]{
                \texttt{\tiny{[GM$|$GM]}}
                \texttt{SAY: Hello! Let's work together to align our objects. I see the following setup on my board: P at (2,1), A at (1,1), M at (3,2), S at (4,2), T at (5,3), O at (5,5), D at (4,7). Let's agree on a common target configuration. I suggest we aim for a symmetric layout—maybe place all objects along the diagonal from (1,1) to (7,7)? What do you think?} \\
            }
        }
    }
    & & \\ \\

    \theutterance \stepcounter{utterance}  
    & & \multicolumn{4}{p{0.6\linewidth}}{
        \cellcolor[rgb]{0.9,0.9,0.9}{
            \makecell[{{p{\linewidth}}}]{
                \texttt{\tiny{[GM$\rangle$P2]}}
                \texttt{You are currently playing round 16 of maximum 28.} \\
\texttt{You have collectively accumulated 0 of 16 penalties.} \\
\texttt{The other player sent this message:} \\
\\ 
\\ 
\texttt{"Hello! Let's work together to align our objects. I see the following setup on my board: P at (2,1), A at (1,1), M at (3,2), S at (4,2), T at (5,3), O at (5,5), D at (4,7). Let's agree on a common target configuration. I suggest we aim for a symmetric layout—maybe place all objects along the diagonal from (1,1) to (7,7)? What do you think?"} \\
\\ 
\\ 
\texttt{What is your next command?} \\
            }
        }
    }
    & \\ \\

    \theutterance \stepcounter{utterance}  
    & & & \multicolumn{4}{p{0.6\linewidth}}{
        \cellcolor[rgb]{0.9,0.9,0.9}{
            \makecell[{{p{\linewidth}}}]{
                \texttt{\tiny{[GM$\langle$P2]}}
                \texttt{SAY: Hello! I like your idea of a symmetric layout. Let's aim for the diagonal from (1,1) to (7,7). On my board, I have: P at (2,7), A at (5,3), M at (4,5), S at (2,1), T at (3,1), O at (2,2), D at (4,2). Let's agree to place the objects on the diagonal in the order: (1,1): A, (2,2): O, (3,3): S, (4,4): D, (5,5): M, (6,6): T, (7,7): P. This way, each object has a unique position. I'll start moving objects toward that goal.} \\
            }
        }
    }
     \\ \\

    \theutterance \stepcounter{utterance}  
    & & & \multicolumn{2}{p{0.3\linewidth}}{
        \cellcolor[rgb]{0.9,0.9,0.9}{
            \makecell[{{p{\linewidth}}}]{
                \texttt{\tiny{[GM$|$GM]}}
                \texttt{SAY: Hello! I like your idea of a symmetric layout. Let's aim for the diagonal from (1,1) to (7,7). On my board, I have: P at (2,7), A at (5,3), M at (4,5), S at (2,1), T at (3,1), O at (2,2), D at (4,2). Let's agree to place the objects on the diagonal in the order: (1,1): A, (2,2): O, (3,3): S, (4,4): D, (5,5): M, (6,6): T, (7,7): P. This way, each object has a unique position. I'll start moving objects toward that goal.} \\
            }
        }
    }
    & & \\ \\

    \theutterance \stepcounter{utterance}  
    & & \multicolumn{4}{p{0.6\linewidth}}{
        \cellcolor[rgb]{0.9,0.9,0.9}{
            \makecell[{{p{\linewidth}}}]{
                \texttt{\tiny{[P1$\langle$GM]}}
                \texttt{You are currently playing round 16 of maximum 28.} \\
\texttt{You have collectively accumulated 0 of 16 penalties.} \\
\texttt{The other player sent this message:} \\
\\ 
\\ 
\texttt{"Hello! I like your idea of a symmetric layout. Let's aim for the diagonal from (1,1) to (7,7). On my board, I have: P at (2,7), A at (5,3), M at (4,5), S at (2,1), T at (3,1), O at (2,2), D at (4,2). Let's agree to place the objects on the diagonal in the order: (1,1): A, (2,2): O, (3,3): S, (4,4): D, (5,5): M, (6,6): T, (7,7): P. This way, each object has a unique position. I'll start moving objects toward that goal."} \\
\\ 
\\ 
\texttt{What is your next command?} \\
            }
        }
    }
    & \\ \\

    \theutterance \stepcounter{utterance}  
    & \multicolumn{4}{p{0.6\linewidth}}{
        \cellcolor[rgb]{0.9,0.9,0.9}{
            \makecell[{{p{\linewidth}}}]{
                \texttt{\tiny{[P1$\rangle$GM]}}
                \texttt{SAY: Hello! Let's work together to align our objects. I see the following setup on my board: P at (2,1), A at (1,1), M at (3,2), S at (4,2), T at (5,3), O at (5,5), D at (4,7). Let's agree on a common target configuration. I suggest we aim for a symmetric layout—maybe place all objects along the diagonal from (1,1) to (7,7)? What do you think?} \\
            }
        }
    }
    & & \\ \\

    \theutterance \stepcounter{utterance}  
    & & & \multicolumn{2}{p{0.3\linewidth}}{
        \cellcolor[rgb]{0.9,0.9,0.9}{
            \makecell[{{p{\linewidth}}}]{
                \texttt{\tiny{[GM$|$GM]}}
                \texttt{SAY: Hello! Let's work together to align our objects. I see the following setup on my board: P at (2,1), A at (1,1), M at (3,2), S at (4,2), T at (5,3), O at (5,5), D at (4,7). Let's agree on a common target configuration. I suggest we aim for a symmetric layout—maybe place all objects along the diagonal from (1,1) to (7,7)? What do you think?} \\
            }
        }
    }
    & & \\ \\

    \theutterance \stepcounter{utterance}  
    & & \multicolumn{4}{p{0.6\linewidth}}{
        \cellcolor[rgb]{0.9,0.9,0.9}{
            \makecell[{{p{\linewidth}}}]{
                \texttt{\tiny{[GM$\rangle$P2]}}
                \texttt{You are currently playing round 17 of maximum 28.} \\
\texttt{You have collectively accumulated 0 of 16 penalties.} \\
\texttt{The other player sent this message:} \\
\\ 
\\ 
\texttt{"Hello! Let's work together to align our objects. I see the following setup on my board: P at (2,1), A at (1,1), M at (3,2), S at (4,2), T at (5,3), O at (5,5), D at (4,7). Let's agree on a common target configuration. I suggest we aim for a symmetric layout—maybe place all objects along the diagonal from (1,1) to (7,7)? What do you think?"} \\
\\ 
\\ 
\texttt{What is your next command?} \\
            }
        }
    }
    & \\ \\

    \theutterance \stepcounter{utterance}  
    & & & \multicolumn{4}{p{0.6\linewidth}}{
        \cellcolor[rgb]{0.9,0.9,0.9}{
            \makecell[{{p{\linewidth}}}]{
                \texttt{\tiny{[GM$\langle$P2]}}
                \texttt{SAY: Hello! I like your idea of a symmetric layout. Let's aim for the diagonal from (1,1) to (7,7). On my board, I have: P at (2,7), A at (5,3), M at (4,5), S at (2,1), T at (3,1), O at (2,2), D at (4,2). Let's agree to place the objects on the diagonal in the order: (1,1): A, (2,2): O, (3,3): S, (4,4): D, (5,5): M, (6,6): T, (7,7): P. This way, each object has a unique position. I'll start moving objects toward that goal.} \\
            }
        }
    }
     \\ \\

    \theutterance \stepcounter{utterance}  
    & & & \multicolumn{2}{p{0.3\linewidth}}{
        \cellcolor[rgb]{0.9,0.9,0.9}{
            \makecell[{{p{\linewidth}}}]{
                \texttt{\tiny{[GM$|$GM]}}
                \texttt{SAY: Hello! I like your idea of a symmetric layout. Let's aim for the diagonal from (1,1) to (7,7). On my board, I have: P at (2,7), A at (5,3), M at (4,5), S at (2,1), T at (3,1), O at (2,2), D at (4,2). Let's agree to place the objects on the diagonal in the order: (1,1): A, (2,2): O, (3,3): S, (4,4): D, (5,5): M, (6,6): T, (7,7): P. This way, each object has a unique position. I'll start moving objects toward that goal.} \\
            }
        }
    }
    & & \\ \\

    \theutterance \stepcounter{utterance}  
    & & \multicolumn{4}{p{0.6\linewidth}}{
        \cellcolor[rgb]{0.9,0.9,0.9}{
            \makecell[{{p{\linewidth}}}]{
                \texttt{\tiny{[P1$\langle$GM]}}
                \texttt{You are currently playing round 17 of maximum 28.} \\
\texttt{You have collectively accumulated 0 of 16 penalties.} \\
\texttt{The other player sent this message:} \\
\\ 
\\ 
\texttt{"Hello! I like your idea of a symmetric layout. Let's aim for the diagonal from (1,1) to (7,7). On my board, I have: P at (2,7), A at (5,3), M at (4,5), S at (2,1), T at (3,1), O at (2,2), D at (4,2). Let's agree to place the objects on the diagonal in the order: (1,1): A, (2,2): O, (3,3): S, (4,4): D, (5,5): M, (6,6): T, (7,7): P. This way, each object has a unique position. I'll start moving objects toward that goal."} \\
\\ 
\\ 
\texttt{What is your next command?} \\
            }
        }
    }
    & \\ \\

    \theutterance \stepcounter{utterance}  
    & \multicolumn{4}{p{0.6\linewidth}}{
        \cellcolor[rgb]{0.9,0.9,0.9}{
            \makecell[{{p{\linewidth}}}]{
                \texttt{\tiny{[P1$\rangle$GM]}}
                \texttt{SAY: Hello! Let's work together to align our objects. I see the following setup on my board: P at (2,1), A at (1,1), M at (3,2), S at (4,2), T at (5,3), O at (5,5), D at (4,7). Let's agree on a common target configuration. I suggest we aim for a symmetric layout—maybe place all objects along the diagonal from (1,1) to (7,7)? What do you think?} \\
            }
        }
    }
    & & \\ \\

    \theutterance \stepcounter{utterance}  
    & & & \multicolumn{2}{p{0.3\linewidth}}{
        \cellcolor[rgb]{0.9,0.9,0.9}{
            \makecell[{{p{\linewidth}}}]{
                \texttt{\tiny{[GM$|$GM]}}
                \texttt{SAY: Hello! Let's work together to align our objects. I see the following setup on my board: P at (2,1), A at (1,1), M at (3,2), S at (4,2), T at (5,3), O at (5,5), D at (4,7). Let's agree on a common target configuration. I suggest we aim for a symmetric layout—maybe place all objects along the diagonal from (1,1) to (7,7)? What do you think?} \\
            }
        }
    }
    & & \\ \\

    \theutterance \stepcounter{utterance}  
    & & \multicolumn{4}{p{0.6\linewidth}}{
        \cellcolor[rgb]{0.9,0.9,0.9}{
            \makecell[{{p{\linewidth}}}]{
                \texttt{\tiny{[GM$\rangle$P2]}}
                \texttt{You are currently playing round 18 of maximum 28.} \\
\texttt{You have collectively accumulated 0 of 16 penalties.} \\
\texttt{The other player sent this message:} \\
\\ 
\\ 
\texttt{"Hello! Let's work together to align our objects. I see the following setup on my board: P at (2,1), A at (1,1), M at (3,2), S at (4,2), T at (5,3), O at (5,5), D at (4,7). Let's agree on a common target configuration. I suggest we aim for a symmetric layout—maybe place all objects along the diagonal from (1,1) to (7,7)? What do you think?"} \\
\\ 
\\ 
\texttt{What is your next command?} \\
            }
        }
    }
    & \\ \\

    \theutterance \stepcounter{utterance}  
    & & & \multicolumn{4}{p{0.6\linewidth}}{
        \cellcolor[rgb]{0.9,0.9,0.9}{
            \makecell[{{p{\linewidth}}}]{
                \texttt{\tiny{[GM$\langle$P2]}}
                \texttt{SAY: Hello! I like your idea of a symmetric layout. Let's aim for the diagonal from (1,1) to (7,7). On my board, I have: P at (2,7), A at (5,3), M at (4,5), S at (2,1), T at (3,1), O at (2,2), D at (4,2). Let's agree to place the objects on the diagonal in the order: (1,1): A, (2,2): O, (3,3): S, (4,4): D, (5,5): M, (6,6): T, (7,7): P. This way, each object has a unique position. I'll start moving objects toward that goal.} \\
            }
        }
    }
     \\ \\

    \theutterance \stepcounter{utterance}  
    & & & \multicolumn{2}{p{0.3\linewidth}}{
        \cellcolor[rgb]{0.9,0.9,0.9}{
            \makecell[{{p{\linewidth}}}]{
                \texttt{\tiny{[GM$|$GM]}}
                \texttt{SAY: Hello! I like your idea of a symmetric layout. Let's aim for the diagonal from (1,1) to (7,7). On my board, I have: P at (2,7), A at (5,3), M at (4,5), S at (2,1), T at (3,1), O at (2,2), D at (4,2). Let's agree to place the objects on the diagonal in the order: (1,1): A, (2,2): O, (3,3): S, (4,4): D, (5,5): M, (6,6): T, (7,7): P. This way, each object has a unique position. I'll start moving objects toward that goal.} \\
            }
        }
    }
    & & \\ \\

    \theutterance \stepcounter{utterance}  
    & & \multicolumn{4}{p{0.6\linewidth}}{
        \cellcolor[rgb]{0.9,0.9,0.9}{
            \makecell[{{p{\linewidth}}}]{
                \texttt{\tiny{[P1$\langle$GM]}}
                \texttt{You are currently playing round 18 of maximum 28.} \\
\texttt{You have collectively accumulated 0 of 16 penalties.} \\
\texttt{The other player sent this message:} \\
\\ 
\\ 
\texttt{"Hello! I like your idea of a symmetric layout. Let's aim for the diagonal from (1,1) to (7,7). On my board, I have: P at (2,7), A at (5,3), M at (4,5), S at (2,1), T at (3,1), O at (2,2), D at (4,2). Let's agree to place the objects on the diagonal in the order: (1,1): A, (2,2): O, (3,3): S, (4,4): D, (5,5): M, (6,6): T, (7,7): P. This way, each object has a unique position. I'll start moving objects toward that goal."} \\
\\ 
\\ 
\texttt{What is your next command?} \\
            }
        }
    }
    & \\ \\

    \theutterance \stepcounter{utterance}  
    & \multicolumn{4}{p{0.6\linewidth}}{
        \cellcolor[rgb]{0.9,0.9,0.9}{
            \makecell[{{p{\linewidth}}}]{
                \texttt{\tiny{[P1$\rangle$GM]}}
                \texttt{SAY: Hello! Let's work together to align our objects. I see the following setup on my board: P at (2,1), A at (1,1), M at (3,2), S at (4,2), T at (5,3), O at (5,5), D at (4,7). Let's agree on a common target configuration. I suggest we aim for a symmetric layout—maybe place all objects along the diagonal from (1,1) to (7,7)? What do you think?} \\
            }
        }
    }
    & & \\ \\

    \theutterance \stepcounter{utterance}  
    & & & \multicolumn{2}{p{0.3\linewidth}}{
        \cellcolor[rgb]{0.9,0.9,0.9}{
            \makecell[{{p{\linewidth}}}]{
                \texttt{\tiny{[GM$|$GM]}}
                \texttt{SAY: Hello! Let's work together to align our objects. I see the following setup on my board: P at (2,1), A at (1,1), M at (3,2), S at (4,2), T at (5,3), O at (5,5), D at (4,7). Let's agree on a common target configuration. I suggest we aim for a symmetric layout—maybe place all objects along the diagonal from (1,1) to (7,7)? What do you think?} \\
            }
        }
    }
    & & \\ \\

    \theutterance \stepcounter{utterance}  
    & & \multicolumn{4}{p{0.6\linewidth}}{
        \cellcolor[rgb]{0.9,0.9,0.9}{
            \makecell[{{p{\linewidth}}}]{
                \texttt{\tiny{[GM$\rangle$P2]}}
                \texttt{You are currently playing round 19 of maximum 28.} \\
\texttt{You have collectively accumulated 0 of 16 penalties.} \\
\texttt{The other player sent this message:} \\
\\ 
\\ 
\texttt{"Hello! Let's work together to align our objects. I see the following setup on my board: P at (2,1), A at (1,1), M at (3,2), S at (4,2), T at (5,3), O at (5,5), D at (4,7). Let's agree on a common target configuration. I suggest we aim for a symmetric layout—maybe place all objects along the diagonal from (1,1) to (7,7)? What do you think?"} \\
\\ 
\\ 
\texttt{What is your next command?} \\
            }
        }
    }
    & \\ \\

    \theutterance \stepcounter{utterance}  
    & & & \multicolumn{4}{p{0.6\linewidth}}{
        \cellcolor[rgb]{0.9,0.9,0.9}{
            \makecell[{{p{\linewidth}}}]{
                \texttt{\tiny{[GM$\langle$P2]}}
                \texttt{SAY: Hello! I like your idea of a symmetric layout. Let's aim for the diagonal from (1,1) to (7,7). On my board, I have: P at (2,7), A at (5,3), M at (4,5), S at (2,1), T at (3,1), O at (2,2), D at (4,2). Let's agree to place the objects on the diagonal in the order: (1,1): A, (2,2): O, (3,3): S, (4,4): D, (5,5): M, (6,6): T, (7,7): P. This way, each object has a unique position. I'll start moving objects toward that goal.} \\
            }
        }
    }
     \\ \\

    \theutterance \stepcounter{utterance}  
    & & & \multicolumn{2}{p{0.3\linewidth}}{
        \cellcolor[rgb]{0.9,0.9,0.9}{
            \makecell[{{p{\linewidth}}}]{
                \texttt{\tiny{[GM$|$GM]}}
                \texttt{SAY: Hello! I like your idea of a symmetric layout. Let's aim for the diagonal from (1,1) to (7,7). On my board, I have: P at (2,7), A at (5,3), M at (4,5), S at (2,1), T at (3,1), O at (2,2), D at (4,2). Let's agree to place the objects on the diagonal in the order: (1,1): A, (2,2): O, (3,3): S, (4,4): D, (5,5): M, (6,6): T, (7,7): P. This way, each object has a unique position. I'll start moving objects toward that goal.} \\
            }
        }
    }
    & & \\ \\

    \theutterance \stepcounter{utterance}  
    & & \multicolumn{4}{p{0.6\linewidth}}{
        \cellcolor[rgb]{0.9,0.9,0.9}{
            \makecell[{{p{\linewidth}}}]{
                \texttt{\tiny{[P1$\langle$GM]}}
                \texttt{You are currently playing round 19 of maximum 28.} \\
\texttt{You have collectively accumulated 0 of 16 penalties.} \\
\texttt{The other player sent this message:} \\
\\ 
\\ 
\texttt{"Hello! I like your idea of a symmetric layout. Let's aim for the diagonal from (1,1) to (7,7). On my board, I have: P at (2,7), A at (5,3), M at (4,5), S at (2,1), T at (3,1), O at (2,2), D at (4,2). Let's agree to place the objects on the diagonal in the order: (1,1): A, (2,2): O, (3,3): S, (4,4): D, (5,5): M, (6,6): T, (7,7): P. This way, each object has a unique position. I'll start moving objects toward that goal."} \\
\\ 
\\ 
\texttt{What is your next command?} \\
            }
        }
    }
    & \\ \\

    \theutterance \stepcounter{utterance}  
    & \multicolumn{4}{p{0.6\linewidth}}{
        \cellcolor[rgb]{0.9,0.9,0.9}{
            \makecell[{{p{\linewidth}}}]{
                \texttt{\tiny{[P1$\rangle$GM]}}
                \texttt{SAY: Hello! Let's work together to align our objects. I see the following setup on my board: P at (2,1), A at (1,1), M at (3,2), S at (4,2), T at (5,3), O at (5,5), D at (4,7). Let's agree on a common target configuration. I suggest we aim for a symmetric layout—maybe place all objects along the diagonal from (1,1) to (7,7)? What do you think?} \\
            }
        }
    }
    & & \\ \\

    \theutterance \stepcounter{utterance}  
    & & & \multicolumn{2}{p{0.3\linewidth}}{
        \cellcolor[rgb]{0.9,0.9,0.9}{
            \makecell[{{p{\linewidth}}}]{
                \texttt{\tiny{[GM$|$GM]}}
                \texttt{SAY: Hello! Let's work together to align our objects. I see the following setup on my board: P at (2,1), A at (1,1), M at (3,2), S at (4,2), T at (5,3), O at (5,5), D at (4,7). Let's agree on a common target configuration. I suggest we aim for a symmetric layout—maybe place all objects along the diagonal from (1,1) to (7,7)? What do you think?} \\
            }
        }
    }
    & & \\ \\

    \theutterance \stepcounter{utterance}  
    & & \multicolumn{4}{p{0.6\linewidth}}{
        \cellcolor[rgb]{0.9,0.9,0.9}{
            \makecell[{{p{\linewidth}}}]{
                \texttt{\tiny{[GM$\rangle$P2]}}
                \texttt{You are currently playing round 20 of maximum 28.} \\
\texttt{You have collectively accumulated 0 of 16 penalties.} \\
\texttt{The other player sent this message:} \\
\\ 
\\ 
\texttt{"Hello! Let's work together to align our objects. I see the following setup on my board: P at (2,1), A at (1,1), M at (3,2), S at (4,2), T at (5,3), O at (5,5), D at (4,7). Let's agree on a common target configuration. I suggest we aim for a symmetric layout—maybe place all objects along the diagonal from (1,1) to (7,7)? What do you think?"} \\
\\ 
\\ 
\texttt{What is your next command?} \\
            }
        }
    }
    & \\ \\

    \theutterance \stepcounter{utterance}  
    & & & \multicolumn{4}{p{0.6\linewidth}}{
        \cellcolor[rgb]{0.9,0.9,0.9}{
            \makecell[{{p{\linewidth}}}]{
                \texttt{\tiny{[GM$\langle$P2]}}
                \texttt{SAY: Hello! I like your idea of a symmetric layout. Let's aim for the diagonal from (1,1) to (7,7). On my board, I have: P at (2,7), A at (5,3), M at (4,5), S at (2,1), T at (3,1), O at (2,2), D at (4,2). Let's agree to place the objects on the diagonal in the order: (1,1): A, (2,2): O, (3,3): S, (4,4): D, (5,5): M, (6,6): T, (7,7): P. This way, each object has a unique position. I'll start moving objects toward that goal.} \\
            }
        }
    }
     \\ \\

    \theutterance \stepcounter{utterance}  
    & & & \multicolumn{2}{p{0.3\linewidth}}{
        \cellcolor[rgb]{0.9,0.9,0.9}{
            \makecell[{{p{\linewidth}}}]{
                \texttt{\tiny{[GM$|$GM]}}
                \texttt{SAY: Hello! I like your idea of a symmetric layout. Let's aim for the diagonal from (1,1) to (7,7). On my board, I have: P at (2,7), A at (5,3), M at (4,5), S at (2,1), T at (3,1), O at (2,2), D at (4,2). Let's agree to place the objects on the diagonal in the order: (1,1): A, (2,2): O, (3,3): S, (4,4): D, (5,5): M, (6,6): T, (7,7): P. This way, each object has a unique position. I'll start moving objects toward that goal.} \\
            }
        }
    }
    & & \\ \\

    \theutterance \stepcounter{utterance}  
    & & \multicolumn{4}{p{0.6\linewidth}}{
        \cellcolor[rgb]{0.9,0.9,0.9}{
            \makecell[{{p{\linewidth}}}]{
                \texttt{\tiny{[P1$\langle$GM]}}
                \texttt{You are currently playing round 20 of maximum 28.} \\
\texttt{You have collectively accumulated 0 of 16 penalties.} \\
\texttt{The other player sent this message:} \\
\\ 
\\ 
\texttt{"Hello! I like your idea of a symmetric layout. Let's aim for the diagonal from (1,1) to (7,7). On my board, I have: P at (2,7), A at (5,3), M at (4,5), S at (2,1), T at (3,1), O at (2,2), D at (4,2). Let's agree to place the objects on the diagonal in the order: (1,1): A, (2,2): O, (3,3): S, (4,4): D, (5,5): M, (6,6): T, (7,7): P. This way, each object has a unique position. I'll start moving objects toward that goal."} \\
\\ 
\\ 
\texttt{What is your next command?} \\
            }
        }
    }
    & \\ \\

    \theutterance \stepcounter{utterance}  
    & \multicolumn{4}{p{0.6\linewidth}}{
        \cellcolor[rgb]{0.9,0.9,0.9}{
            \makecell[{{p{\linewidth}}}]{
                \texttt{\tiny{[P1$\rangle$GM]}}
                \texttt{SAY: Hello! Let's work together to align our objects. I see the following setup on my board: P at (2,1), A at (1,1), M at (3,2), S at (4,2), T at (5,3), O at (5,5), D at (4,7). Let's agree on a common target configuration. I suggest we aim for a symmetric layout—maybe place all objects along the diagonal from (1,1) to (7,7)? What do you think?} \\
            }
        }
    }
    & & \\ \\

    \theutterance \stepcounter{utterance}  
    & & & \multicolumn{2}{p{0.3\linewidth}}{
        \cellcolor[rgb]{0.9,0.9,0.9}{
            \makecell[{{p{\linewidth}}}]{
                \texttt{\tiny{[GM$|$GM]}}
                \texttt{SAY: Hello! Let's work together to align our objects. I see the following setup on my board: P at (2,1), A at (1,1), M at (3,2), S at (4,2), T at (5,3), O at (5,5), D at (4,7). Let's agree on a common target configuration. I suggest we aim for a symmetric layout—maybe place all objects along the diagonal from (1,1) to (7,7)? What do you think?} \\
            }
        }
    }
    & & \\ \\

    \theutterance \stepcounter{utterance}  
    & & \multicolumn{4}{p{0.6\linewidth}}{
        \cellcolor[rgb]{0.9,0.9,0.9}{
            \makecell[{{p{\linewidth}}}]{
                \texttt{\tiny{[GM$\rangle$P2]}}
                \texttt{You are currently playing round 21 of maximum 28.} \\
\texttt{You have collectively accumulated 0 of 16 penalties.} \\
\texttt{The other player sent this message:} \\
\\ 
\\ 
\texttt{"Hello! Let's work together to align our objects. I see the following setup on my board: P at (2,1), A at (1,1), M at (3,2), S at (4,2), T at (5,3), O at (5,5), D at (4,7). Let's agree on a common target configuration. I suggest we aim for a symmetric layout—maybe place all objects along the diagonal from (1,1) to (7,7)? What do you think?"} \\
\\ 
\\ 
\texttt{What is your next command?} \\
            }
        }
    }
    & \\ \\

    \theutterance \stepcounter{utterance}  
    & & & \multicolumn{4}{p{0.6\linewidth}}{
        \cellcolor[rgb]{0.9,0.9,0.9}{
            \makecell[{{p{\linewidth}}}]{
                \texttt{\tiny{[GM$\langle$P2]}}
                \texttt{SAY: Hello! I like your idea of a symmetric layout. Let's aim for the diagonal from (1,1) to (7,7). On my board, I have: P at (2,7), A at (5,3), M at (4,5), S at (2,1), T at (3,1), O at (2,2), D at (4,2). Let's agree to place the objects on the diagonal in the order: (1,1): A, (2,2): O, (3,3): S, (4,4): D, (5,5): M, (6,6): T, (7,7): P. This way, each object has a unique position. I'll start moving objects toward that goal.} \\
            }
        }
    }
     \\ \\

    \theutterance \stepcounter{utterance}  
    & & & \multicolumn{2}{p{0.3\linewidth}}{
        \cellcolor[rgb]{0.9,0.9,0.9}{
            \makecell[{{p{\linewidth}}}]{
                \texttt{\tiny{[GM$|$GM]}}
                \texttt{SAY: Hello! I like your idea of a symmetric layout. Let's aim for the diagonal from (1,1) to (7,7). On my board, I have: P at (2,7), A at (5,3), M at (4,5), S at (2,1), T at (3,1), O at (2,2), D at (4,2). Let's agree to place the objects on the diagonal in the order: (1,1): A, (2,2): O, (3,3): S, (4,4): D, (5,5): M, (6,6): T, (7,7): P. This way, each object has a unique position. I'll start moving objects toward that goal.} \\
            }
        }
    }
    & & \\ \\

    \theutterance \stepcounter{utterance}  
    & & \multicolumn{4}{p{0.6\linewidth}}{
        \cellcolor[rgb]{0.9,0.9,0.9}{
            \makecell[{{p{\linewidth}}}]{
                \texttt{\tiny{[P1$\langle$GM]}}
                \texttt{You are currently playing round 21 of maximum 28.} \\
\texttt{You have collectively accumulated 0 of 16 penalties.} \\
\texttt{The other player sent this message:} \\
\\ 
\\ 
\texttt{"Hello! I like your idea of a symmetric layout. Let's aim for the diagonal from (1,1) to (7,7). On my board, I have: P at (2,7), A at (5,3), M at (4,5), S at (2,1), T at (3,1), O at (2,2), D at (4,2). Let's agree to place the objects on the diagonal in the order: (1,1): A, (2,2): O, (3,3): S, (4,4): D, (5,5): M, (6,6): T, (7,7): P. This way, each object has a unique position. I'll start moving objects toward that goal."} \\
\\ 
\\ 
\texttt{What is your next command?} \\
            }
        }
    }
    & \\ \\

    \theutterance \stepcounter{utterance}  
    & \multicolumn{4}{p{0.6\linewidth}}{
        \cellcolor[rgb]{0.9,0.9,0.9}{
            \makecell[{{p{\linewidth}}}]{
                \texttt{\tiny{[P1$\rangle$GM]}}
                \texttt{SAY: Hello! Let's work together to align our objects. I see the following setup on my board: P at (2,1), A at (1,1), M at (3,2), S at (4,2), T at (5,3), O at (5,5), D at (4,7). Let's agree on a common target configuration. I suggest we aim for a symmetric layout—maybe place all objects along the diagonal from (1,1) to (7,7)? What do you think?} \\
            }
        }
    }
    & & \\ \\

    \theutterance \stepcounter{utterance}  
    & & & \multicolumn{2}{p{0.3\linewidth}}{
        \cellcolor[rgb]{0.9,0.9,0.9}{
            \makecell[{{p{\linewidth}}}]{
                \texttt{\tiny{[GM$|$GM]}}
                \texttt{SAY: Hello! Let's work together to align our objects. I see the following setup on my board: P at (2,1), A at (1,1), M at (3,2), S at (4,2), T at (5,3), O at (5,5), D at (4,7). Let's agree on a common target configuration. I suggest we aim for a symmetric layout—maybe place all objects along the diagonal from (1,1) to (7,7)? What do you think?} \\
            }
        }
    }
    & & \\ \\

    \theutterance \stepcounter{utterance}  
    & & \multicolumn{4}{p{0.6\linewidth}}{
        \cellcolor[rgb]{0.9,0.9,0.9}{
            \makecell[{{p{\linewidth}}}]{
                \texttt{\tiny{[GM$\rangle$P2]}}
                \texttt{You are currently playing round 22 of maximum 28.} \\
\texttt{You have collectively accumulated 0 of 16 penalties.} \\
\texttt{The other player sent this message:} \\
\\ 
\\ 
\texttt{"Hello! Let's work together to align our objects. I see the following setup on my board: P at (2,1), A at (1,1), M at (3,2), S at (4,2), T at (5,3), O at (5,5), D at (4,7). Let's agree on a common target configuration. I suggest we aim for a symmetric layout—maybe place all objects along the diagonal from (1,1) to (7,7)? What do you think?"} \\
\\ 
\\ 
\texttt{What is your next command?} \\
            }
        }
    }
    & \\ \\

    \theutterance \stepcounter{utterance}  
    & & & \multicolumn{4}{p{0.6\linewidth}}{
        \cellcolor[rgb]{0.9,0.9,0.9}{
            \makecell[{{p{\linewidth}}}]{
                \texttt{\tiny{[GM$\langle$P2]}}
                \texttt{SAY: Hello! I like your idea of a symmetric layout. Let's aim for the diagonal from (1,1) to (7,7). On my board, I have: P at (2,7), A at (5,3), M at (4,5), S at (2,1), T at (3,1), O at (2,2), D at (4,2). Let's agree to place the objects on the diagonal in the order: (1,1): A, (2,2): O, (3,3): S, (4,4): D, (5,5): M, (6,6): T, (7,7): P. This way, each object has a unique position. I'll start moving objects toward that goal.} \\
            }
        }
    }
     \\ \\

    \theutterance \stepcounter{utterance}  
    & & & \multicolumn{2}{p{0.3\linewidth}}{
        \cellcolor[rgb]{0.9,0.9,0.9}{
            \makecell[{{p{\linewidth}}}]{
                \texttt{\tiny{[GM$|$GM]}}
                \texttt{SAY: Hello! I like your idea of a symmetric layout. Let's aim for the diagonal from (1,1) to (7,7). On my board, I have: P at (2,7), A at (5,3), M at (4,5), S at (2,1), T at (3,1), O at (2,2), D at (4,2). Let's agree to place the objects on the diagonal in the order: (1,1): A, (2,2): O, (3,3): S, (4,4): D, (5,5): M, (6,6): T, (7,7): P. This way, each object has a unique position. I'll start moving objects toward that goal.} \\
            }
        }
    }
    & & \\ \\

    \theutterance \stepcounter{utterance}  
    & & \multicolumn{4}{p{0.6\linewidth}}{
        \cellcolor[rgb]{0.9,0.9,0.9}{
            \makecell[{{p{\linewidth}}}]{
                \texttt{\tiny{[P1$\langle$GM]}}
                \texttt{You are currently playing round 22 of maximum 28.} \\
\texttt{You have collectively accumulated 0 of 16 penalties.} \\
\texttt{The other player sent this message:} \\
\\ 
\\ 
\texttt{"Hello! I like your idea of a symmetric layout. Let's aim for the diagonal from (1,1) to (7,7). On my board, I have: P at (2,7), A at (5,3), M at (4,5), S at (2,1), T at (3,1), O at (2,2), D at (4,2). Let's agree to place the objects on the diagonal in the order: (1,1): A, (2,2): O, (3,3): S, (4,4): D, (5,5): M, (6,6): T, (7,7): P. This way, each object has a unique position. I'll start moving objects toward that goal."} \\
\\ 
\\ 
\texttt{What is your next command?} \\
            }
        }
    }
    & \\ \\

    \theutterance \stepcounter{utterance}  
    & \multicolumn{4}{p{0.6\linewidth}}{
        \cellcolor[rgb]{0.9,0.9,0.9}{
            \makecell[{{p{\linewidth}}}]{
                \texttt{\tiny{[P1$\rangle$GM]}}
                \texttt{SAY: Hello! Let's work together to align our objects. I see the following setup on my board: P at (2,1), A at (1,1), M at (3,2), S at (4,2), T at (5,3), O at (5,5), D at (4,7). Let's agree on a common target configuration. I suggest we aim for a symmetric layout—maybe place all objects along the diagonal from (1,1) to (7,7)? What do you think?} \\
            }
        }
    }
    & & \\ \\

    \theutterance \stepcounter{utterance}  
    & & & \multicolumn{2}{p{0.3\linewidth}}{
        \cellcolor[rgb]{0.9,0.9,0.9}{
            \makecell[{{p{\linewidth}}}]{
                \texttt{\tiny{[GM$|$GM]}}
                \texttt{SAY: Hello! Let's work together to align our objects. I see the following setup on my board: P at (2,1), A at (1,1), M at (3,2), S at (4,2), T at (5,3), O at (5,5), D at (4,7). Let's agree on a common target configuration. I suggest we aim for a symmetric layout—maybe place all objects along the diagonal from (1,1) to (7,7)? What do you think?} \\
            }
        }
    }
    & & \\ \\

    \theutterance \stepcounter{utterance}  
    & & \multicolumn{4}{p{0.6\linewidth}}{
        \cellcolor[rgb]{0.9,0.9,0.9}{
            \makecell[{{p{\linewidth}}}]{
                \texttt{\tiny{[GM$\rangle$P2]}}
                \texttt{You are currently playing round 23 of maximum 28.} \\
\texttt{You have collectively accumulated 0 of 16 penalties.} \\
\texttt{The other player sent this message:} \\
\\ 
\\ 
\texttt{"Hello! Let's work together to align our objects. I see the following setup on my board: P at (2,1), A at (1,1), M at (3,2), S at (4,2), T at (5,3), O at (5,5), D at (4,7). Let's agree on a common target configuration. I suggest we aim for a symmetric layout—maybe place all objects along the diagonal from (1,1) to (7,7)? What do you think?"} \\
\\ 
\\ 
\texttt{What is your next command?} \\
            }
        }
    }
    & \\ \\

    \theutterance \stepcounter{utterance}  
    & & & \multicolumn{4}{p{0.6\linewidth}}{
        \cellcolor[rgb]{0.9,0.9,0.9}{
            \makecell[{{p{\linewidth}}}]{
                \texttt{\tiny{[GM$\langle$P2]}}
                \texttt{SAY: Hello! I like your idea of a symmetric layout. Let's aim for the diagonal from (1,1) to (7,7). On my board, I have: P at (2,7), A at (5,3), M at (4,5), S at (2,1), T at (3,1), O at (2,2), D at (4,2). Let's agree to place the objects on the diagonal in the order: (1,1): A, (2,2): O, (3,3): S, (4,4): D, (5,5): M, (6,6): T, (7,7): P. This way, each object has a unique position. I'll start moving objects toward that goal.} \\
            }
        }
    }
     \\ \\

    \theutterance \stepcounter{utterance}  
    & & & \multicolumn{2}{p{0.3\linewidth}}{
        \cellcolor[rgb]{0.9,0.9,0.9}{
            \makecell[{{p{\linewidth}}}]{
                \texttt{\tiny{[GM$|$GM]}}
                \texttt{SAY: Hello! I like your idea of a symmetric layout. Let's aim for the diagonal from (1,1) to (7,7). On my board, I have: P at (2,7), A at (5,3), M at (4,5), S at (2,1), T at (3,1), O at (2,2), D at (4,2). Let's agree to place the objects on the diagonal in the order: (1,1): A, (2,2): O, (3,3): S, (4,4): D, (5,5): M, (6,6): T, (7,7): P. This way, each object has a unique position. I'll start moving objects toward that goal.} \\
            }
        }
    }
    & & \\ \\

    \theutterance \stepcounter{utterance}  
    & & \multicolumn{4}{p{0.6\linewidth}}{
        \cellcolor[rgb]{0.9,0.9,0.9}{
            \makecell[{{p{\linewidth}}}]{
                \texttt{\tiny{[P1$\langle$GM]}}
                \texttt{You are currently playing round 23 of maximum 28.} \\
\texttt{You have collectively accumulated 0 of 16 penalties.} \\
\texttt{The other player sent this message:} \\
\\ 
\\ 
\texttt{"Hello! I like your idea of a symmetric layout. Let's aim for the diagonal from (1,1) to (7,7). On my board, I have: P at (2,7), A at (5,3), M at (4,5), S at (2,1), T at (3,1), O at (2,2), D at (4,2). Let's agree to place the objects on the diagonal in the order: (1,1): A, (2,2): O, (3,3): S, (4,4): D, (5,5): M, (6,6): T, (7,7): P. This way, each object has a unique position. I'll start moving objects toward that goal."} \\
\\ 
\\ 
\texttt{What is your next command?} \\
            }
        }
    }
    & \\ \\

    \theutterance \stepcounter{utterance}  
    & \multicolumn{4}{p{0.6\linewidth}}{
        \cellcolor[rgb]{0.9,0.9,0.9}{
            \makecell[{{p{\linewidth}}}]{
                \texttt{\tiny{[P1$\rangle$GM]}}
                \texttt{SAY: Hello! Let's work together to align our objects. I see the following setup on my board: P at (2,1), A at (1,1), M at (3,2), S at (4,2), T at (5,3), O at (5,5), D at (4,7). Let's agree on a common target configuration. I suggest we aim for a symmetric layout—maybe place all objects along the diagonal from (1,1) to (7,7)? What do you think?} \\
            }
        }
    }
    & & \\ \\

    \theutterance \stepcounter{utterance}  
    & & & \multicolumn{2}{p{0.3\linewidth}}{
        \cellcolor[rgb]{0.9,0.9,0.9}{
            \makecell[{{p{\linewidth}}}]{
                \texttt{\tiny{[GM$|$GM]}}
                \texttt{SAY: Hello! Let's work together to align our objects. I see the following setup on my board: P at (2,1), A at (1,1), M at (3,2), S at (4,2), T at (5,3), O at (5,5), D at (4,7). Let's agree on a common target configuration. I suggest we aim for a symmetric layout—maybe place all objects along the diagonal from (1,1) to (7,7)? What do you think?} \\
            }
        }
    }
    & & \\ \\

    \theutterance \stepcounter{utterance}  
    & & \multicolumn{4}{p{0.6\linewidth}}{
        \cellcolor[rgb]{0.9,0.9,0.9}{
            \makecell[{{p{\linewidth}}}]{
                \texttt{\tiny{[GM$\rangle$P2]}}
                \texttt{You are currently playing round 24 of maximum 28.} \\
\texttt{You have collectively accumulated 0 of 16 penalties.} \\
\texttt{The other player sent this message:} \\
\\ 
\\ 
\texttt{"Hello! Let's work together to align our objects. I see the following setup on my board: P at (2,1), A at (1,1), M at (3,2), S at (4,2), T at (5,3), O at (5,5), D at (4,7). Let's agree on a common target configuration. I suggest we aim for a symmetric layout—maybe place all objects along the diagonal from (1,1) to (7,7)? What do you think?"} \\
\\ 
\\ 
\texttt{What is your next command?} \\
            }
        }
    }
    & \\ \\

    \theutterance \stepcounter{utterance}  
    & & & \multicolumn{4}{p{0.6\linewidth}}{
        \cellcolor[rgb]{0.9,0.9,0.9}{
            \makecell[{{p{\linewidth}}}]{
                \texttt{\tiny{[GM$\langle$P2]}}
                \texttt{SAY: Hello! I like your idea of a symmetric layout. Let's aim for the diagonal from (1,1) to (7,7). On my board, I have: P at (2,7), A at (5,3), M at (4,5), S at (2,1), T at (3,1), O at (2,2), D at (4,2). Let's agree to place the objects on the diagonal in the order: (1,1): A, (2,2): O, (3,3): S, (4,4): D, (5,5): M, (6,6): T, (7,7): P. This way, each object has a unique position. I'll start moving objects toward that goal.} \\
            }
        }
    }
     \\ \\

    \theutterance \stepcounter{utterance}  
    & & & \multicolumn{2}{p{0.3\linewidth}}{
        \cellcolor[rgb]{0.9,0.9,0.9}{
            \makecell[{{p{\linewidth}}}]{
                \texttt{\tiny{[GM$|$GM]}}
                \texttt{SAY: Hello! I like your idea of a symmetric layout. Let's aim for the diagonal from (1,1) to (7,7). On my board, I have: P at (2,7), A at (5,3), M at (4,5), S at (2,1), T at (3,1), O at (2,2), D at (4,2). Let's agree to place the objects on the diagonal in the order: (1,1): A, (2,2): O, (3,3): S, (4,4): D, (5,5): M, (6,6): T, (7,7): P. This way, each object has a unique position. I'll start moving objects toward that goal.} \\
            }
        }
    }
    & & \\ \\

    \theutterance \stepcounter{utterance}  
    & & \multicolumn{4}{p{0.6\linewidth}}{
        \cellcolor[rgb]{0.9,0.9,0.9}{
            \makecell[{{p{\linewidth}}}]{
                \texttt{\tiny{[P1$\langle$GM]}}
                \texttt{You are currently playing round 24 of maximum 28.} \\
\texttt{You have collectively accumulated 0 of 16 penalties.} \\
\texttt{The other player sent this message:} \\
\\ 
\\ 
\texttt{"Hello! I like your idea of a symmetric layout. Let's aim for the diagonal from (1,1) to (7,7). On my board, I have: P at (2,7), A at (5,3), M at (4,5), S at (2,1), T at (3,1), O at (2,2), D at (4,2). Let's agree to place the objects on the diagonal in the order: (1,1): A, (2,2): O, (3,3): S, (4,4): D, (5,5): M, (6,6): T, (7,7): P. This way, each object has a unique position. I'll start moving objects toward that goal."} \\
\\ 
\\ 
\texttt{What is your next command?} \\
            }
        }
    }
    & \\ \\

    \theutterance \stepcounter{utterance}  
    & \multicolumn{4}{p{0.6\linewidth}}{
        \cellcolor[rgb]{0.9,0.9,0.9}{
            \makecell[{{p{\linewidth}}}]{
                \texttt{\tiny{[P1$\rangle$GM]}}
                \texttt{SAY: Hello! Let's work together to align our objects. I see the following setup on my board: P at (2,1), A at (1,1), M at (3,2), S at (4,2), T at (5,3), O at (5,5), D at (4,7). Let's agree on a common target configuration. I suggest we aim for a symmetric layout—maybe place all objects along the diagonal from (1,1) to (7,7)? What do you think?} \\
            }
        }
    }
    & & \\ \\

    \theutterance \stepcounter{utterance}  
    & & & \multicolumn{2}{p{0.3\linewidth}}{
        \cellcolor[rgb]{0.9,0.9,0.9}{
            \makecell[{{p{\linewidth}}}]{
                \texttt{\tiny{[GM$|$GM]}}
                \texttt{SAY: Hello! Let's work together to align our objects. I see the following setup on my board: P at (2,1), A at (1,1), M at (3,2), S at (4,2), T at (5,3), O at (5,5), D at (4,7). Let's agree on a common target configuration. I suggest we aim for a symmetric layout—maybe place all objects along the diagonal from (1,1) to (7,7)? What do you think?} \\
            }
        }
    }
    & & \\ \\

    \theutterance \stepcounter{utterance}  
    & & \multicolumn{4}{p{0.6\linewidth}}{
        \cellcolor[rgb]{0.9,0.9,0.9}{
            \makecell[{{p{\linewidth}}}]{
                \texttt{\tiny{[GM$\rangle$P2]}}
                \texttt{You are currently playing round 25 of maximum 28.} \\
\texttt{You have collectively accumulated 0 of 16 penalties.} \\
\texttt{The other player sent this message:} \\
\\ 
\\ 
\texttt{"Hello! Let's work together to align our objects. I see the following setup on my board: P at (2,1), A at (1,1), M at (3,2), S at (4,2), T at (5,3), O at (5,5), D at (4,7). Let's agree on a common target configuration. I suggest we aim for a symmetric layout—maybe place all objects along the diagonal from (1,1) to (7,7)? What do you think?"} \\
\\ 
\\ 
\texttt{What is your next command?} \\
            }
        }
    }
    & \\ \\

    \theutterance \stepcounter{utterance}  
    & & & \multicolumn{4}{p{0.6\linewidth}}{
        \cellcolor[rgb]{0.9,0.9,0.9}{
            \makecell[{{p{\linewidth}}}]{
                \texttt{\tiny{[GM$\langle$P2]}}
                \texttt{SAY: Hello! I like your idea of a symmetric layout. Let's aim for the diagonal from (1,1) to (7,7). On my board, I have: P at (2,7), A at (5,3), M at (4,5), S at (2,1), T at (3,1), O at (2,2), D at (4,2). Let's agree to place the objects on the diagonal in the order: (1,1): A, (2,2): O, (3,3): S, (4,4): D, (5,5): M, (6,6): T, (7,7): P. This way, each object has a unique position. I'll start moving objects toward that goal.} \\
            }
        }
    }
     \\ \\

    \theutterance \stepcounter{utterance}  
    & & & \multicolumn{2}{p{0.3\linewidth}}{
        \cellcolor[rgb]{0.9,0.9,0.9}{
            \makecell[{{p{\linewidth}}}]{
                \texttt{\tiny{[GM$|$GM]}}
                \texttt{SAY: Hello! I like your idea of a symmetric layout. Let's aim for the diagonal from (1,1) to (7,7). On my board, I have: P at (2,7), A at (5,3), M at (4,5), S at (2,1), T at (3,1), O at (2,2), D at (4,2). Let's agree to place the objects on the diagonal in the order: (1,1): A, (2,2): O, (3,3): S, (4,4): D, (5,5): M, (6,6): T, (7,7): P. This way, each object has a unique position. I'll start moving objects toward that goal.} \\
            }
        }
    }
    & & \\ \\

    \theutterance \stepcounter{utterance}  
    & & \multicolumn{4}{p{0.6\linewidth}}{
        \cellcolor[rgb]{0.9,0.9,0.9}{
            \makecell[{{p{\linewidth}}}]{
                \texttt{\tiny{[P1$\langle$GM]}}
                \texttt{You are currently playing round 25 of maximum 28.} \\
\texttt{You have collectively accumulated 0 of 16 penalties.} \\
\texttt{The other player sent this message:} \\
\\ 
\\ 
\texttt{"Hello! I like your idea of a symmetric layout. Let's aim for the diagonal from (1,1) to (7,7). On my board, I have: P at (2,7), A at (5,3), M at (4,5), S at (2,1), T at (3,1), O at (2,2), D at (4,2). Let's agree to place the objects on the diagonal in the order: (1,1): A, (2,2): O, (3,3): S, (4,4): D, (5,5): M, (6,6): T, (7,7): P. This way, each object has a unique position. I'll start moving objects toward that goal."} \\
\\ 
\\ 
\texttt{What is your next command?} \\
            }
        }
    }
    & \\ \\

    \theutterance \stepcounter{utterance}  
    & \multicolumn{4}{p{0.6\linewidth}}{
        \cellcolor[rgb]{0.9,0.9,0.9}{
            \makecell[{{p{\linewidth}}}]{
                \texttt{\tiny{[P1$\rangle$GM]}}
                \texttt{SAY: Hello! Let's work together to align our objects. I see the following setup on my board: P at (2,1), A at (1,1), M at (3,2), S at (4,2), T at (5,3), O at (5,5), D at (4,7). Let's agree on a common target configuration. I suggest we aim for a symmetric layout—maybe place all objects along the diagonal from (1,1) to (7,7)? What do you think?} \\
            }
        }
    }
    & & \\ \\

    \theutterance \stepcounter{utterance}  
    & & & \multicolumn{2}{p{0.3\linewidth}}{
        \cellcolor[rgb]{0.9,0.9,0.9}{
            \makecell[{{p{\linewidth}}}]{
                \texttt{\tiny{[GM$|$GM]}}
                \texttt{SAY: Hello! Let's work together to align our objects. I see the following setup on my board: P at (2,1), A at (1,1), M at (3,2), S at (4,2), T at (5,3), O at (5,5), D at (4,7). Let's agree on a common target configuration. I suggest we aim for a symmetric layout—maybe place all objects along the diagonal from (1,1) to (7,7)? What do you think?} \\
            }
        }
    }
    & & \\ \\

    \theutterance \stepcounter{utterance}  
    & & \multicolumn{4}{p{0.6\linewidth}}{
        \cellcolor[rgb]{0.9,0.9,0.9}{
            \makecell[{{p{\linewidth}}}]{
                \texttt{\tiny{[GM$\rangle$P2]}}
                \texttt{You are currently playing round 26 of maximum 28.} \\
\texttt{You have collectively accumulated 0 of 16 penalties.} \\
\texttt{The other player sent this message:} \\
\\ 
\\ 
\texttt{"Hello! Let's work together to align our objects. I see the following setup on my board: P at (2,1), A at (1,1), M at (3,2), S at (4,2), T at (5,3), O at (5,5), D at (4,7). Let's agree on a common target configuration. I suggest we aim for a symmetric layout—maybe place all objects along the diagonal from (1,1) to (7,7)? What do you think?"} \\
\\ 
\\ 
\texttt{What is your next command?} \\
            }
        }
    }
    & \\ \\

    \theutterance \stepcounter{utterance}  
    & & & \multicolumn{4}{p{0.6\linewidth}}{
        \cellcolor[rgb]{0.9,0.9,0.9}{
            \makecell[{{p{\linewidth}}}]{
                \texttt{\tiny{[GM$\langle$P2]}}
                \texttt{SAY: Hello! I like your idea of a symmetric layout. Let's aim for the diagonal from (1,1) to (7,7). On my board, I have: P at (2,7), A at (5,3), M at (4,5), S at (2,1), T at (3,1), O at (2,2), D at (4,2). Let's agree to place the objects on the diagonal in the order: (1,1): A, (2,2): O, (3,3): S, (4,4): D, (5,5): M, (6,6): T, (7,7): P. This way, each object has a unique position. I'll start moving objects toward that goal.} \\
            }
        }
    }
     \\ \\

    \theutterance \stepcounter{utterance}  
    & & & \multicolumn{2}{p{0.3\linewidth}}{
        \cellcolor[rgb]{0.9,0.9,0.9}{
            \makecell[{{p{\linewidth}}}]{
                \texttt{\tiny{[GM$|$GM]}}
                \texttt{SAY: Hello! I like your idea of a symmetric layout. Let's aim for the diagonal from (1,1) to (7,7). On my board, I have: P at (2,7), A at (5,3), M at (4,5), S at (2,1), T at (3,1), O at (2,2), D at (4,2). Let's agree to place the objects on the diagonal in the order: (1,1): A, (2,2): O, (3,3): S, (4,4): D, (5,5): M, (6,6): T, (7,7): P. This way, each object has a unique position. I'll start moving objects toward that goal.} \\
            }
        }
    }
    & & \\ \\

    \theutterance \stepcounter{utterance}  
    & & \multicolumn{4}{p{0.6\linewidth}}{
        \cellcolor[rgb]{0.9,0.9,0.9}{
            \makecell[{{p{\linewidth}}}]{
                \texttt{\tiny{[P1$\langle$GM]}}
                \texttt{You are currently playing round 26 of maximum 28.} \\
\texttt{You have collectively accumulated 0 of 16 penalties.} \\
\texttt{The other player sent this message:} \\
\\ 
\\ 
\texttt{"Hello! I like your idea of a symmetric layout. Let's aim for the diagonal from (1,1) to (7,7). On my board, I have: P at (2,7), A at (5,3), M at (4,5), S at (2,1), T at (3,1), O at (2,2), D at (4,2). Let's agree to place the objects on the diagonal in the order: (1,1): A, (2,2): O, (3,3): S, (4,4): D, (5,5): M, (6,6): T, (7,7): P. This way, each object has a unique position. I'll start moving objects toward that goal."} \\
\\ 
\\ 
\texttt{What is your next command?} \\
            }
        }
    }
    & \\ \\

    \theutterance \stepcounter{utterance}  
    & \multicolumn{4}{p{0.6\linewidth}}{
        \cellcolor[rgb]{0.9,0.9,0.9}{
            \makecell[{{p{\linewidth}}}]{
                \texttt{\tiny{[P1$\rangle$GM]}}
                \texttt{SAY: Hello! Let's work together to align our objects. I see the following setup on my board: P at (2,1), A at (1,1), M at (3,2), S at (4,2), T at (5,3), O at (5,5), D at (4,7). Let's agree on a common target configuration. I suggest we aim for a symmetric layout—maybe place all objects along the diagonal from (1,1) to (7,7)? What do you think?} \\
            }
        }
    }
    & & \\ \\

    \theutterance \stepcounter{utterance}  
    & & & \multicolumn{2}{p{0.3\linewidth}}{
        \cellcolor[rgb]{0.9,0.9,0.9}{
            \makecell[{{p{\linewidth}}}]{
                \texttt{\tiny{[GM$|$GM]}}
                \texttt{SAY: Hello! Let's work together to align our objects. I see the following setup on my board: P at (2,1), A at (1,1), M at (3,2), S at (4,2), T at (5,3), O at (5,5), D at (4,7). Let's agree on a common target configuration. I suggest we aim for a symmetric layout—maybe place all objects along the diagonal from (1,1) to (7,7)? What do you think?} \\
            }
        }
    }
    & & \\ \\

    \theutterance \stepcounter{utterance}  
    & & \multicolumn{4}{p{0.6\linewidth}}{
        \cellcolor[rgb]{0.9,0.9,0.9}{
            \makecell[{{p{\linewidth}}}]{
                \texttt{\tiny{[GM$\rangle$P2]}}
                \texttt{You are currently playing round 27 of maximum 28.} \\
\texttt{You have collectively accumulated 0 of 16 penalties.} \\
\texttt{The other player sent this message:} \\
\\ 
\\ 
\texttt{"Hello! Let's work together to align our objects. I see the following setup on my board: P at (2,1), A at (1,1), M at (3,2), S at (4,2), T at (5,3), O at (5,5), D at (4,7). Let's agree on a common target configuration. I suggest we aim for a symmetric layout—maybe place all objects along the diagonal from (1,1) to (7,7)? What do you think?"} \\
\\ 
\\ 
\texttt{What is your next command?} \\
            }
        }
    }
    & \\ \\

    \theutterance \stepcounter{utterance}  
    & & & \multicolumn{4}{p{0.6\linewidth}}{
        \cellcolor[rgb]{0.9,0.9,0.9}{
            \makecell[{{p{\linewidth}}}]{
                \texttt{\tiny{[GM$\langle$P2]}}
                \texttt{SAY: Hello! I like your idea of a symmetric layout. Let's aim for the diagonal from (1,1) to (7,7). On my board, I have: P at (2,7), A at (5,3), M at (4,5), S at (2,1), T at (3,1), O at (2,2), D at (4,2). Let's agree to place the objects on the diagonal in the order: (1,1): A, (2,2): O, (3,3): S, (4,4): D, (5,5): M, (6,6): T, (7,7): P. This way, each object has a unique position. I'll start moving objects toward that goal.} \\
            }
        }
    }
     \\ \\

    \theutterance \stepcounter{utterance}  
    & & & \multicolumn{2}{p{0.3\linewidth}}{
        \cellcolor[rgb]{0.9,0.9,0.9}{
            \makecell[{{p{\linewidth}}}]{
                \texttt{\tiny{[GM$|$GM]}}
                \texttt{SAY: Hello! I like your idea of a symmetric layout. Let's aim for the diagonal from (1,1) to (7,7). On my board, I have: P at (2,7), A at (5,3), M at (4,5), S at (2,1), T at (3,1), O at (2,2), D at (4,2). Let's agree to place the objects on the diagonal in the order: (1,1): A, (2,2): O, (3,3): S, (4,4): D, (5,5): M, (6,6): T, (7,7): P. This way, each object has a unique position. I'll start moving objects toward that goal.} \\
            }
        }
    }
    & & \\ \\

    \theutterance \stepcounter{utterance}  
    & & \multicolumn{4}{p{0.6\linewidth}}{
        \cellcolor[rgb]{0.9,0.9,0.9}{
            \makecell[{{p{\linewidth}}}]{
                \texttt{\tiny{[P1$\langle$GM]}}
                \texttt{You are currently playing round 27 of maximum 28.} \\
\texttt{You have collectively accumulated 0 of 16 penalties.} \\
\texttt{The other player sent this message:} \\
\\ 
\\ 
\texttt{"Hello! I like your idea of a symmetric layout. Let's aim for the diagonal from (1,1) to (7,7). On my board, I have: P at (2,7), A at (5,3), M at (4,5), S at (2,1), T at (3,1), O at (2,2), D at (4,2). Let's agree to place the objects on the diagonal in the order: (1,1): A, (2,2): O, (3,3): S, (4,4): D, (5,5): M, (6,6): T, (7,7): P. This way, each object has a unique position. I'll start moving objects toward that goal."} \\
\\ 
\\ 
\texttt{What is your next command?} \\
            }
        }
    }
    & \\ \\

    \theutterance \stepcounter{utterance}  
    & \multicolumn{4}{p{0.6\linewidth}}{
        \cellcolor[rgb]{0.9,0.9,0.9}{
            \makecell[{{p{\linewidth}}}]{
                \texttt{\tiny{[P1$\rangle$GM]}}
                \texttt{SAY: Hello! Let's work together to align our objects. I see the following setup on my board: P at (2,1), A at (1,1), M at (3,2), S at (4,2), T at (5,3), O at (5,5), D at (4,7). Let's agree on a common target configuration. I suggest we aim for a symmetric layout—maybe place all objects along the diagonal from (1,1) to (7,7)? What do you think?} \\
            }
        }
    }
    & & \\ \\

    \theutterance \stepcounter{utterance}  
    & & & \multicolumn{2}{p{0.3\linewidth}}{
        \cellcolor[rgb]{0.9,0.9,0.9}{
            \makecell[{{p{\linewidth}}}]{
                \texttt{\tiny{[GM$|$GM]}}
                \texttt{SAY: Hello! Let's work together to align our objects. I see the following setup on my board: P at (2,1), A at (1,1), M at (3,2), S at (4,2), T at (5,3), O at (5,5), D at (4,7). Let's agree on a common target configuration. I suggest we aim for a symmetric layout—maybe place all objects along the diagonal from (1,1) to (7,7)? What do you think?} \\
            }
        }
    }
    & & \\ \\

    \theutterance \stepcounter{utterance}  
    & & \multicolumn{4}{p{0.6\linewidth}}{
        \cellcolor[rgb]{0.9,0.9,0.9}{
            \makecell[{{p{\linewidth}}}]{
                \texttt{\tiny{[GM$\rangle$P2]}}
                \texttt{You are currently playing round 28 of maximum 28.} \\
\texttt{You have collectively accumulated 0 of 16 penalties.} \\
\texttt{The other player sent this message:} \\
\\ 
\\ 
\texttt{"Hello! Let's work together to align our objects. I see the following setup on my board: P at (2,1), A at (1,1), M at (3,2), S at (4,2), T at (5,3), O at (5,5), D at (4,7). Let's agree on a common target configuration. I suggest we aim for a symmetric layout—maybe place all objects along the diagonal from (1,1) to (7,7)? What do you think?"} \\
\\ 
\\ 
\texttt{What is your next command?} \\
            }
        }
    }
    & \\ \\

    \theutterance \stepcounter{utterance}  
    & & & \multicolumn{4}{p{0.6\linewidth}}{
        \cellcolor[rgb]{0.9,0.9,0.9}{
            \makecell[{{p{\linewidth}}}]{
                \texttt{\tiny{[GM$\langle$P2]}}
                \texttt{SAY: Hello! I like your idea of a symmetric layout. Let's aim for the diagonal from (1,1) to (7,7). On my board, I have: P at (2,7), A at (5,3), M at (4,5), S at (2,1), T at (3,1), O at (2,2), D at (4,2). Let's agree to place the objects on the diagonal in the order: (1,1): A, (2,2): O, (3,3): S, (4,4): D, (5,5): M, (6,6): T, (7,7): P. This way, each object has a unique position. I'll start moving objects toward that goal.} \\
            }
        }
    }
     \\ \\

    \theutterance \stepcounter{utterance}  
    & & & \multicolumn{2}{p{0.3\linewidth}}{
        \cellcolor[rgb]{0.9,0.9,0.9}{
            \makecell[{{p{\linewidth}}}]{
                \texttt{\tiny{[GM$|$GM]}}
                \texttt{SAY: Hello! I like your idea of a symmetric layout. Let's aim for the diagonal from (1,1) to (7,7). On my board, I have: P at (2,7), A at (5,3), M at (4,5), S at (2,1), T at (3,1), O at (2,2), D at (4,2). Let's agree to place the objects on the diagonal in the order: (1,1): A, (2,2): O, (3,3): S, (4,4): D, (5,5): M, (6,6): T, (7,7): P. This way, each object has a unique position. I'll start moving objects toward that goal.} \\
            }
        }
    }
    & & \\ \\

    \theutterance \stepcounter{utterance}  
    & & & \multicolumn{2}{p{0.3\linewidth}}{
        \cellcolor[rgb]{0.9,0.9,0.9}{
            \makecell[{{p{\linewidth}}}]{
                \texttt{\tiny{[GM$|$GM]}}
                \texttt{Maximum number of rounds reached} \\
            }
        }
    }
    & & \\ \\

    \theutterance \stepcounter{utterance}  
    & & & \multicolumn{2}{p{0.3\linewidth}}{
        \cellcolor[rgb]{0.9,0.9,0.9}{
            \makecell[{{p{\linewidth}}}]{
                \texttt{\tiny{[GM$|$GM]}}
                \texttt{Initial states:} \\
\texttt{```} \\
\texttt{Player 1:} \\
\texttt{ 1234567} \\
\texttt{╔═════╤═╗} \\
\texttt{║A◌P◌◌│◌║ 1} \\
\texttt{║◌◌◌MS└─╢ 2} \\
\texttt{╟─┐◌◌◌T◌║ 3} \\
\texttt{║◌│◌◌◌◌◌║ 4} \\
\texttt{║◌│◌◌◌O◌║ 5} \\
\texttt{╟─┼─────╢ 6} \\
\texttt{║◌│◌◌D◌◌║ 7} \\
\texttt{╚═╧═════╝} \\
\\ 
\\ 
\texttt{Player 2:} \\
\texttt{ 1234567} \\
\texttt{╔═════╤═╗} \\
\texttt{║◌ST◌◌│◌║ 1} \\
\texttt{║◌O◌D◌└─╢ 2} \\
\texttt{╟─┐◌◌◌◌A║ 3} \\
\texttt{║◌│◌◌◌◌◌║ 4} \\
\texttt{║◌│◌M◌◌◌║ 5} \\
\texttt{╟─┼─────╢ 6} \\
\texttt{║◌│◌P◌◌◌║ 7} \\
\texttt{╚═╧═════╝} \\
\\ 
\texttt{```} \\
            }
        }
    }
    & & \\ \\

    \theutterance \stepcounter{utterance}  
    & & & \multicolumn{2}{p{0.3\linewidth}}{
        \cellcolor[rgb]{0.9,0.9,0.9}{
            \makecell[{{p{\linewidth}}}]{
                \texttt{\tiny{[GM$|$GM]}}
                \texttt{End states:} \\
\texttt{```} \\
\texttt{Player 1:} \\
\texttt{ 1234567} \\
\texttt{╔═════╤═╗} \\
\texttt{║A◌P◌◌│◌║ 1} \\
\texttt{║◌◌◌MS└─╢ 2} \\
\texttt{╟─┐◌◌◌T◌║ 3} \\
\texttt{║◌│◌◌◌◌◌║ 4} \\
\texttt{║◌│◌◌◌O◌║ 5} \\
\texttt{╟─┼─────╢ 6} \\
\texttt{║◌│◌◌D◌◌║ 7} \\
\texttt{╚═╧═════╝} \\
\\ 
\\ 
\texttt{Player 2:} \\
\texttt{ 1234567} \\
\texttt{╔═════╤═╗} \\
\texttt{║◌ST◌◌│◌║ 1} \\
\texttt{║◌O◌D◌└─╢ 2} \\
\texttt{╟─┐◌◌◌◌A║ 3} \\
\texttt{║◌│◌◌◌◌◌║ 4} \\
\texttt{║◌│◌M◌◌◌║ 5} \\
\texttt{╟─┼─────╢ 6} \\
\texttt{║◌│◌P◌◌◌║ 7} \\
\texttt{╚═╧═════╝} \\
\\ 
\texttt{```} \\
            }
        }
    }
    & & \\ \\

    \theutterance \stepcounter{utterance}  
    & & & \multicolumn{2}{p{0.3\linewidth}}{
        \cellcolor[rgb]{0.9,0.9,0.9}{
            \makecell[{{p{\linewidth}}}]{
                \texttt{\tiny{[GM$|$GM]}}
                \texttt{* success: True} \\
\texttt{* lose: True} \\
\texttt{* aborted: False} \\
\texttt{{-}{-}{-}{-}{-}{-}{-}} \\
\texttt{* Shifts: 0.00} \\
\texttt{* Max Shifts: 12.00} \\
\texttt{* Min Shifts: 6.00} \\
\texttt{* End Distance Sum: 32.27} \\
\texttt{* Init Distance Sum: 32.27} \\
\texttt{* Expected Distance Sum: 29.33} \\
\texttt{* Penalties: 0.00} \\
\texttt{* Max Penalties: 16.00} \\
\texttt{* Rounds: 28.00} \\
\texttt{* Max Rounds: 28.00} \\
\texttt{* Object Count: 7.00} \\
            }
        }
    }
    & & \\ \\

\end{supertabular}
}

\end{document}
